\begin{chapterbox}
    \chapter{Der Giftzwerg und die Sphäre (2020)}
    \label{Der Giftzwerg und die Sphäre (2020)}
    \az{0 bis 340}
    
    Im Anschluss an die Düsteren Zeiten sitzt ein Mechaniker in einem Kerker der Schildzwerge fest. Nun bleibt ihm bloß sein Verstand, um die Geschehnisse in Cavern (und einen neugierigen Fürsten) zu seinen Gunsten zu lenken. Kann Kjall seiner Zelle entkommen, den Lauf der Zeit ins Lot rücken und in seine Heimat zurückkehren? Will er das überhaupt?
\end{chapterbox}


\section{Prolog: Bragor verscherbelt den Bruderschild}

\az{Jahr 62}

Das Rietgras wogte sachte im Wind. Eine massige Gestalt wanderte über das Feld. Nur mit einem Lendenschurz aus Schaffell und großen Stiefeln war sie bekleidet. Auf ihrem Kopf saßen zwei mächtige Hörner. In ihrer einen Hand hielt sie einen Speer. Ihre andere Hand war über kunstvolle Knoten mit einem beschmutzten Schild verbunden, den die Gestalt hinter sich her schleifte. Das Sonnenlicht ließ die prallen Muskeln des Wesens glänzen. Es war ein Tarus.

Eine kleine Gestalt trat aus dem Schatten eines Baumes hervor und näherte sich dem Tarus. Sie trug einen großen Rucksack und benutzte einen schweren Schmiedehammer als Stock. In der Hand hielt sie eine Art silberner Kugel. Ein roter Haarschopf umrahmte ein kantiges Gesicht. Es war ein Zwerg.

„Seid herzlich gegrüßt!“, rief der Tarus mit tiefer Stimme fröhlich und hob dann seine Hand, an welcher der dreckige Schild festgeknotet war, „Ihr hättet nicht zufälligerweise Interesse daran, einen Schild zu kaufen?“

Der Zwerg lächelte, zeigte gelbe Zähne und sprach: „Zufälligerweise hätte ich Interesse daran.“

Der Tarus grinste breit zurück: „Wie viel wäre dieser Schild wert, was meint Ihr? Genug, damit man damit etwas Ausrüstung für eine Reise in ein Gebirge erwerben könnte?“

„Da bin ich mir nicht sicher. Wenn ich mir den so ansehe, würde ich nicht mehr als ein Goldstück dafür springen lassen.“

„Was kann ich denn für ein Goldstück erhalten?“, fragte der Tarus neugierig.

Der Zwerg stutzte kurz und schmunzelte dann: „Viel, äußerst viel. Genug für mehrere Wochen Proviant und die beste Kletterausrüstung östlich des Fahlen Gebirges. Nicht, dass Ihr letzteres nötig hättet, wenn ich Euch so ansehe.“

Der Zwerg holte einen kleinen Beutel hervor, in dem einige Münzen klirrten, und zog eine Goldmünze heraus. Sie war klein und die Prägung darauf saß schief, doch der Tarus betrachtete sie dennoch fasziniert. Dass ein schäbiger Schild gleich viel wert sein sollte wie ganze Säcke voller Apfelnüsse und eine wertvolle Kletterausrüstung, schien ihm nicht seltsam vorzukommen. Und so wurde der Tausch schnell abgehandelt.

Tarus und Zwerg waren kurz davor, ihrer eigenen Wege zu gehen, da hielt der Tarus nochmals inne und fragte: „Meister Zwerg, wisst ihr vielleicht, wo eine Pflanze namens ‚Sternkraut‘ wächst?“

Der Zwerg musterte den Tarus lange und meinte dann: „Ihr erinnert mich an jemanden, den ich einst kannte... sagt, wie lautet Euer Name?“

„Bragor, Meister Zwerg. Und eurer?“

„Ich sehe schon überall Geister...“, murmelte der Zwerg mehr zu sich selbst als zum Tarus, „Mein Name ist für Euch nicht von Belang. Und nein, werter Bragor, ich kenne mich leider nicht in der Kunde der Pflanzen aus. Ich würde Euch raten, zu einer Kräuterhexe zu gehen, oder vielleicht zum Baum der Lieder, da wird man Euch eher aushelfen können.“

Der Tarus ließ seine Schultern hängen: „Beim Baum der Lieder war ich schon“, murmelte er.

Der Zwerg kümmerte sich nicht weiter darum, sondern betrachtete den dreckigen Schild in seinen Händen ehrfürchtig. Dann drehte er sich um, winkte dem Tarus zum Abschied zu und lief dem Sonnenuntergang dagegen. Er war in Hochstimmung. Er hatte soeben den zweiten der vier mächtigen Schilde aus uralter Zeit errungen.








\newpage
\section{Gefangen}

\az{Jahr 0}



Kjall saß in seiner Kerkerzelle und schrieb. Die Schildzwerge waren so zuvorkommend gewesen, ihm ein wackeliges Pult, durchnässtes Pergament und vertrocknete Tinte zur Verfügung zu stellen, damit er sich den lieben langen Tag dort verwirklichen könnte. Natürlich war sich Kjall bewusst, dass die Wachen alle paar Tage seine beschriebenen Pergamente durchstöberten, um zu schauen, ob sie darin irgendwelche Hinweise auf seine Herkunft oder seine Tiefminen-Golems finden könnten. Dachten sie tatsächlich, dass er es nicht mitbekommen würde, wenn sie mitten in der Nacht laut polternd an sein Pult traten und in seinen Schriften rumwühlten?

Kjall hörte dennoch nicht mit Schreiben und Skizzieren auf. Außer dem Pult und seiner Pritsche befand sich in dieser Zelle kaum etwas, auch keine weiteren Vergnügungsmöglichkeiten. Ausgang hatte er nebst seiner halbtäglichen Toilettengänge auch kaum. Er musste nur darauf achten, keine wichtigen Informationen niederzuschreiben. Und so verbrachte Kjall die meiste Zeit damit, Schmiergedichte auf seine Wachen zu verfassen sowie Skizzen kreativer Konstruktionen, die abbildeten, was Kjall mit den Wachen gerne anstellen würde – natürlich gespickt mit furiosen Flüchen verschiedenster Art. Es war niederer Humor, aber dadurch nicht minder amüsant.\bigskip



Schwere Schritte hallten durch den dunklen Gang, an dessen Ende Kjalls Kerker lag. Es war unüblich für die Wachen, bereits zu dieser Zeit zurückzukehren. Kjalls Magen hatte noch nicht einmal zu knurren begonnen, was heißen musste, dass die nächste Mahlzeit noch lange nicht an der Reihe war. Wollten sie ihn etwa wieder befragen? Inzwischen musste doch selbst der einfältigste dieser Narren begriffen haben, dass er ihnen nicht verraten würde, von wo er kam. Erst recht nicht, nachdem sie die Überreste seiner Tiefminen-Golems eingeschmolzen hatten. Die Golems aus wertvollem Roteisen waren vor der Sphäre sein größter Schatz gewesen, und wenn er hier erst einmal raus war, würden sie es noch bereuen, diese Meisterwerke nicht als solche geachtet zu haben.

Die Schritte im Gang verstummten. Kjall setzte seine Feder ab und drehte sich zu den Besuchern um. Er erwartete, den üblichen missmutig gelaunten Weißbart zu erblicken. Stattdessen fiel sein Blick auf goldene Stiefel und einen verfilzten grauen Bart. Buschige Augenbrauen, die wirre, eisblaue Augen verdeckten. Gesichtszüge, die ihm nur allzu bekannt erschienen.

„Fürst Hallwort!“, hauchte Kjall. Hatte dieser neugierige Trottel sich tatsächlich dazu herabgelassen, einen einfachen Gefangenen wie Kjall zu besuchen, in der Hoffnung, mehr über seine Herkunft und sein Können zu erfahren? Kjall war natürlich nicht bereit, dem Fürsten das Geheimnis der Golemkerne zu verraten. Aber das musste er Hallwort nicht direkt auf die Nase binden. Schließlich besaß der Fürst nicht nur die Schlüssel zu seiner Zelle, sondern auch die Autorität, Kjall laufen zu lassen. Moment mal... Schlüssel! Kjalls Gedanken rasten. Fürst Hallwort besaß nicht nur die Schlüssel zu Kjalls Zelle, sondern auch den zur geheimen Fürstenkammer! Der geheimen Fürstenkammer, in welcher Kjall und seine Kumpanen in einigen Jahrzehnten die Sphäre finden würden, jenes Artefakt, welches Kjall überhaupt erst in die Vergangenheit versetzt hatte. Und das bedeutete, dass die jüngere Version der Sphäre vielleicht in diesem Augenblick in der geheimen Fürstenkammer lag. So nah! Wenn Kjall an diese Sphäre kommen könnte, wäre er nicht nur frei, sondern auch mächtiger als je zuvor. Das würde es ihm erlauben, noch weiter zurück in die Vergangenheit zu reisen, um die Gründung Andors zu verhindern und die Werte der Schildzwerge zu bewahren! Und dafür musste er nur eines hinkriegen: Den neugierigen Idiot vor sich geschickt zu manipulieren.

„Ich sehe mit Freude, dass du meinen Namen kennst“, plapperte Fürst Hallwort los, „Deine Notizen sind in modernen Zwergenrunen geschrieben. Deine Hände sind die vernarbten Hände eines Zwergs, der sein Leben in den Tiefminen verbracht hat. Und doch scheint dich keiner dort zu vermissen. Du bist ein Mysterium, o Zwerg, und du wolltest den Wächtern nicht einmal deinen Namen nennen. Nun frage ich dich, wie ein gewisser Zwergenfürst dich umzustimmen vermögen könnte. Ich spüre, dass an dir mehr steckt, als man auf den ersten Blick meinen mag.“

Kjall traute seinen Ohren kaum. Hatte Hallwort tatsächlich in den Tiefminen umherfragen lassen, ob ihn jemand erkannte? Hatten er und seine Tiefminen-Golems einen derart starken Eindruck hinterlassen? Ein Glück, dass sie nicht sein jüngeres Ich aufgespürt hatten, sonst hätte er einiges zu erklären gehabt.

„Mein Name ist... Jork, mein Fürst“, log Kjall, „Und wenn Ihr mir schon die Ehre erweist, mich hier zu besuchen, so will ich Euch ein Geheimnis verraten...“

Kjall senkte seine Stimme und Hallwort trat näher an sein Gitter. Kjall betrachtete die beiden Zwergenwachen, die auf weiter hinten im Schatten des Gangs warteten und die Interaktion beobachteten. Sie mochten nicht verstehen, was Kjall flüsterte, aber sie würden definitiv mitkriegen, wenn Kjall eine unüberlegte Bewegung tätigte. So blieb Kjall ruhig und wisperte weiter: „Mein Fürst, ich komme aus der Zukunft! Ich bin hier auf der Mission, Cavern vor dem Untergang zu bewahren!“

Gespannt wartete Kjall auf Hallworts Reaktion. Der Fürst ließ sich nichts anmerken, sondern starrte Kjall mit seinem eisblauen Blick geradewegs in die Augen, still darauf wartend, dass Kjall weitersprach. Er war nicht überzeugt, natürlich nicht, aber er hatte ihn auch noch nicht ausgelacht. Hallworts hatte bekanntlich ein offenes Ohr für jede Theorie, die man ihm auf die Nase binden wollte. Und das wusste Kjall zu seinem Vorteil zu nutzen. Jetzt musste er nur noch Argumente finden, die seine Behauptung untermauern würden.

„Mein Fürst, ich bin hier in Eurem Auftrag. Ich hatte das Ziel, mit den mechanischen Golems ins Rietland zu ziehen und einen gefährlichen Andori von einer finsteren Tat abzuhalten.“

„Ach“, sagte Hallwort. Er schwieg eine Zeit lang und streichelte seinen grauen Rauschebart.

„Hmm“, murmelte Hallwort dann, „Äußerst interessant. Aus welchem Jahr meint Ihr zu kommen, Jork? Was ist das für ein gefährlicher Andori und inwiefern würde er Caverns Untergang herbeiführen?“

Kjall sprach langsam und bedacht weiter, während sein Kopf ratterte: „Vielleicht ist es besser, wenn ich Euch so wenig wie möglich üblich die Zukunft verrate.“

Hallwort lachte auf: „Sehr geschickt! Das macht es äußerst schwer, zu überprüfen, ob Ihr die Wahrheit sprecht. Sagt, werter Jork, warum habt Ihr nie nach mir gerufen? Ihr schienet bereit, bis ans Ende Eurer Tage in dieser Zelle zu verstauben. Wenn Ihr tatsächlich in meinem Auftrag hierher geschickt worden wärt, so hättet Ihr doch sicherlich nach einer Audienz bei mir verlangt.“

Kjall biss sich auf die Unterlippe. Er war noch nie sonderlich gut im Improvisieren gewesen, erst recht nicht, wenn so viel auf dem Spiel stand. Wie lautete seine Geschichte? Er war von Fürst Hallwort aus einer Zukunft geschickt worden, in welcher Cavern dem Untergang geweiht war. Hallwort hatte ihn hierher geschickt. Warum hätte er sich still verhalten? Weil, weil...

„Ich kenne Euch doch, mein Fürst! Nichts weckt Euer Interesse so schnell wie ein ungelöstes Rätsel. Ihr seid schneller zu mir gekommen, als wenn ich, ein einfacher Gefangener, um eine Audienz gebeten hätte. Feuer und Spucke, dass Ihr hier vor mir steht, zeigt doch gerade, dass es funktioniert hat! Ich... ich weiß Dinge, die ich nicht wissen könnte, wenn ich löge! Ihr trinkt Euren Drachentee mit immer mit zwei Löffeln Honig drin, auch wenn Ihr das bei jeder passenden Gelegenheit abstreitet. Ihr führt eine geheime Tölplingzucht im Brauneisenstein. Als Kind seid Ihr einmal in den Geheimen See gefallen, als ihr heimlich und unerlaubt Cavern zu erkunden versuchtet, und eine von Wasser triefende, grünliche Gestalt rettete euch.“

Kjall dankte dem feurigen Gott für die Klatschtanten, die sich nach Hallworts Tod über ihn das Maul zerrissen hatten. So vieles war damals aufgedeckt worden, das die breite Bevölkerung Caverns zum jetzigen Zeitpunkt noch gar nicht wissen konnte. Hallwort legte seinen Kopf schief und schien zum ersten Mal ernsthaft zu bedenken, ob Kjall vielleicht die Wahrheit sagte. Nach einer weiteren Überlegungspause fragte er kleinlaut: „Werden meine Tölplinge entdeckt werden?“

Jetzt hab‘ ich dich, dachte Kjall, und sprach: „Arpachen-Überfall. Fragt nicht weiter nach. Niemand wird ernsthaft zu Schaden kommen.“

Stille.

„Ist es Euch denn gelungen, diesen gefährlichen Andori zu überwinden, ehe Ihr gefangen genommen wurdet?“

„Ich befürchte nicht, mein Fürst. Er hat weitere Andori auf mich gehetzt. Schade, dass Ihr meine Golems eingeschmolzen habt. Aber wir werden auch so eine Möglichkeit finden, an ihn heranzukommen.“

Hallwort tippte sich nachdenklich auf die Nase: „Selbst wenn Ihr mich überzeugen könnt, dass Ihr aus der Zukunft und in meinem eigenen Auftrag hier seid, so dürfte es noch einiges schwerer werden, den guten König Brandur davon zu überzeugen. Zu verlangen, dass er einen seiner Landleute herausrücke für etwaige Verbrechen, die noch gar nicht begangen wurden, ist kein geschicktes Vorgehen, um den Frieden zwischen unseren Völkern zu wahren. Und Frieden ist in diesen kriegerischen Zeiten das höchste Gut. Die Trolle...“

„Mein Fürst, deswegen versuchen wir ja auch gar nicht, Brandur davon zu überzeugen. Der fragliche Andori hat keine starke Beziehung zum Königshaus und ich habe in der Vergangenheit keinerlei Beziehung zu euch. Ich muss nur hier raus und kann dann ganz im Geheimen...“

Kjall riss seine Augen auf, als wäre ihm soeben etwas eingefallen: „Mein Fürst, die Apparatur, mit welcher ich in die Vergangenheit reisen konnte, stammt aus Cavern, aus der geheimen Fürstenkammer! Mit etwas Glück befindet sich das Artefakt in diesem Augenblick bereits dort! Wenn ich dessen Funktionsweise demonstrieren könnte, so würde Euch das den letzten Zweifel nehmen, dass ich tatsächlich aus der Zukunft komme.“

„...und da nur ich Zugang zur Kammer habe, wäre damit geklärt, dass ich Euch hierher gesandt habe. Genial!“, rief Hallwort aus, „Schnell, werter Jork, verratet mir, wonach ich Ausschau halten soll!“

Hatte er die Lüge tatsächlich abgekauft? Spielte der Fürst seine Begeisterung nur? Wollte er Kjall auf die Probe stellen? Kjall war sich nicht sicher, aber es konnte sicherlich nicht schaden, mitzuspielen.

„Das Artefakt ist eine metallene Sphäre aus einer Vielzahl ineinander verschlungener Ranken und Verstrebungen. Ihr werdet es wissen, wenn Ihr es seht. Findet diese Sphäre und bringt sie hierher, dann kann ich Euch ihre Funktionsweise demonstrieren.“

Hallwort sah Kjall prüfend in die Augen. Dann nickte er, drehte sich energetisch um und stapfte von Kjalls Zelle weg. Die beiden Zwergenwachen musterten Kjall finster und schritten dann im Einklang ihrem Fürsten hinterher. Kjall blieb alleine zurück und überlegte.\bigskip



Überraschenderweise wurde Kjalls Mediation nicht durch zurückkehrende Wachen oder gar einen neugierigen Fürsten gestört. Stattdessen glomm plötzlich ein blauer Funke in der Mitte der Zelle auf. Noch während Kjall sich aufrichtete, um den Funken genauer zu betrachten, erstrahlte jener und wurde zu einer unförmigen, wabernden blauen Masse. Die Masse wuchs langsam, bis sie beinahe die Größe eines Skrals hatte. Kjall hatte ein solches Phänomen bereits einmal gesehen: Als er und seine Komplizen das Portal in die Vergangenheit geöffnet hatten, das sie hierhergebracht hatte.

So war Kjall eher fasziniert als überrascht, als die blaue Masse zu rotieren begann, immer breiter und dünner wurde, bis sie als irrsinnig rasch entlang ihrer Symmetrieachse rotierende Scheibe in der Mitte seiner Zelle schwebte. Durch diese hellblaue Hülle konnte Kjall verschwommen erkennen, dass auf der anderen Seite des Portals drei Gestalten zu sehen waren, drei Zwerge. Ihre Proportionen waren verzerrt, als würde Kjall sie durch eine dicke Wasserschicht betrachten. Sie schienen sich zu unterhalten.

Kjall streckte seine Hand nach dem blauen Portal aus und spürte, wie dessen Hülle unter seinen Fingerspitzen nachgab. Auf der anderen Seite des Portals schien einer der anderen Zwerge es ihm gleichzutun. Schnell trat Kjall zurück, um Platz zu lassen, denn er antizipierte, was als nächsten geschehen könnte.

Tatsächlich traten die drei Zwerge kurz darauf synchron nach vorne. Blaue Lichtstrahlen brachen aus der Portalhülle hervor und ließen Schatten über die karge Zellenwand tanzen. Mit zugekniffenen Augen versuchte Kjall, die Portaloberfläche im Blick zu halten. Die Schemen der drei Zwerge wogen und waberten, zerfaserten und setzten sich wieder zusammen. Dann brachen sie aus der Oberfläche des Portals hervor und plumpsten mit einem lauten Knall auf den Zellenboden. Das blaue Leuchten ebbte ab und das Portal wurde wieder zur einer blassen hellblauen Scheibe, die in der Mitte des Raums schwebte. Doch Kjall achtete nicht auf sie. Was da vor ihm lag, war um einiges spannender. Der erste der drei Zwerge, die soeben durch das Portal getreten – oder besser gefallen – waren, war eine von Kjalls Wachen. Wie lautete sein Name schon wieder? Bort? Der zweite Zwerg war schon um einiges spannender, denn es handelte sich dabei um Fürst Hallwort. Und der dritte Zwerg... war Kjall selbst!\bigskip



Eine Zeitreise durch ein solches Portal war ein ganz und gar unangenehmes Geschäft, daran mochte sich Kjall noch erinnern. Man schwebte eine Zeit lang durch eine Art Limbo, einem endlosen Raum aus gleißendem blauen Licht, in dem es keine Schwerkraft zum Orientieren und keine Luft zum Atmen gab. Kjall hatte bei seiner ersten Reise gedacht, dass sein letztes Stündlein geschlagen hätte, als er nach rund einer Minute sein Bewusstsein verloren hatte. Stattdessen war er in der Vergangenheit aufgewacht, in den Armen eines finsteren Schattenwesens, welches ihn durchs Rietland trug. Den beiden verfluchten Schatten, dieser Shan und diesem Bleichen, hatte die luftlose Reise durch den Limbo offenbar nichts ausgemacht. Immerhin war Kjall damals nicht der einzige gewesen, der ohnmächtig geworden war, und der mächtige Tarus Thogger war noch um einiges schwerer zu tragen gewesen als er selbst. Dennoch hatte der Vorfall etwas an seinem Stolz gekratzt.

Genauso wie Kjall damals selbst ohnmächtig geworden war, lagen die drei zeitreisenden Zwerge nun allesamt ohnmächtig in einer Reihe vor ihm: Bort, Hallwort und ein zweiter Kjall. Es war ein wahrlich unwirkliches Gefühl, eine Kopie seiner selbst zu betrachten. Sah Kjall wirklich so gedrungen aus? Und was war mit diesen abstehenden Ohren los? Die Hände und Füße des zweiten Kjalls waren mit schweren Eisenketten verbunden, welche wiederum an Borts Gürtel befestigt waren. Hallwort hingegen war genauso ketten- wie bewusstlos. An seinem Gürtel hängte das übliche Sammelsurium an Täschchen und Ketten voller wertlosem Allerlei, wie man es von Fürst Hallwort erwarten würde – aber auch ein kleines, unscheinbares Stück Metall mit einigen Scharnieren und Drähten dran. Ein Casamatuc, ein Zwergenwerkzeug, mit dem man fast jedes Schloss aufkriegen konnte! Was für ein Geschenk des Schicksals!

Kjall fiel kaum auf, dass das durchscheinende Zeitportal seine Form verlor, wieder zu einem kleinen Funken wurde und dann vollends verpuffte. Stattdessen stürzte er sich auf den Gürtel des ohnmächtigen Hallwort, griff nach dem Casamatuc, rannte zur Zellentür und setzte das Werkzeug an. Schnell, schnell, schnell! Wie lange blieb ein kräftiger Zwerg üblicherweise bewusstlos? Kjalls geschickte Finger waren wie geschaffen für solche Tätigkeiten und der Casamatuc des Fürsten des Schildzwerge war wie erwartet von allerbester Qualität. Ein Klick, zwei Klicks, beim dritten klemmte es, etwas Rütteln, ein Klack, eine Drehung, und die Zellentür war offen. Kjall war frei!

Ein kurzer Kontrollblick zurück verriet Kjall, dass Hallwort die Fürstenkrone mit den vier eingemeißelten Schilden aus uralter Zeit zum Treffen mit seinem mysteriösen Gefangenen nicht aufgesetzt hatte. Zu schade, die hätte sich in Kjalls Museum wirklich prächtig gemacht. Aber das war nun wirklich nicht die Hauptsache. Kjall wollte sich bereits wieder abwenden und auf die Socken machen, da besann er sich eines Besseren. Er wusste zwar nicht, aus welcher Zeitlinie dieser zweite Kjall stammte, aber dessen Hallwort war gleich gekleidet wie derjenige Hallwort, welcher sich erst vor Kurzem noch mit Kjall unterhalten hatte. Es bestand bestimmt eine gewisse Chance, dass Kjall selbst bald in eine ähnliche Lage wie dieser zweite Kjall kommen würde, und dann würde er es bestimmt sehr schätzen, nicht angekettet an eine andere Wache zu Bewusstsein zu kommen. Und ohnehin sollte man doch immer Solidarität mit sich selbst zeigen.

Rasch huschte Kjall zurück zum zweiten Kjall, sorgsam darauf achtend, ob Hallwort oder Bort Zeichen des Erwachsens zeigten. Sie taten es nicht, und so kniete sich Kjall neben dem zweiten Kjall nieder und setzte den Casamatuc an dessen Fesseln. Er drehte und stieß, es klickte und knirschte, aber die Ketten blieben zu.

„Bei Bailas angekokelter Eiche!“, knurrte Kjall. Er hätte nicht gedacht, dass es solche Casamatuc-resistenten Schlösser bereits für einfache Eisenketten gab. Was nun? Sollte er den zweiten Kjall doch einfach zurücklassen? Was würde geschehen, wenn man die beiden anderen Zwerge hier finden würde? Was wäre das für eine Zeitlinie, in welcher fortan zwei Borts, und noch schlimmer, zwei Hallworts existierten?

In diesem Moment riss der zweite Kjall seine Augen auf und packte Kjall am Handgelenk. Er richtete sich geschmeidig auf – seine Ketten klirrten laut – und sah sich in der Zelle um.

„Schnell“, richte er dann das Wort an Kjall, „Hilf mir, die beiden anderen aus dieser Zelle zu schaffen. Keine Sorge, das war ihre erste Zeitreise. So schnell werden sie nicht wieder erwachen.“ Bei den Flammen des Zor, klang Kjalls Stimme immer so krächzend?

Ohne auf eine Antwort zu warten, packte der zweite Kjall Bort an den Schultern und begann, die Wache in Richtung der immer noch offen stehenden Zellentür zu bugsieren. Kjall starrte seinem Doppelgänger mit einem flauen Gefühl im Magen nach, fasste sich, beugte sich dann zum ohnmächtigen Fürst Hallwort hinunter und zerrte diesen aus der Zelle heraus. Was hatte dieser Zwerg nur gegessen, schweres Roteisen?! Es war absurd. Der Fürst der Schildzwerge in seinen Händen, ihm komplett ausgeliefert. Und doch konnte er nichts damit anfangen, da es sich um den Hallwort aus einer anderen Zeitlinie handelte!

Kjall folgte dem zweiten Kjall ächzend aus der Zelle hinaus in einen kleinen Seitengang, wo er Fürst Hallwort erleichtert zu Boden plumpsen ließ. Der zweite Kjall hatte Bort dort bereits auf eine leichte Decke gebettet und meinte: „Wir müssen noch mal schnell in die Zelle zurück.“

Die Kette, die die Hände und Füße des zweiten Kjall an Bort knüpfte, war so lang, dass der zweite Kjall bis fast in die Zelle zurücklaufen konnte, ehe seine Kette anschlug.

„Der Casamatuc!“, zischte der zweite Kjall, „Rasch!“

Kjall huschte an ihm vorbei in die Zelle zurück, griff nach dem dort liegenden Casamatuc und übergab das Werkzeug seinem Doppelgänger: „Er wird die Ketten leider nicht lösen können, das habe ich bereits versucht.“

Der zweite Kjall nickte: „Ich ebenfalls.“

Ehe Kjall über diese Worte nachsinnieren konnte, setzte sein Doppelgänger nach: „Aber wir wollen doch keine Spuren zurücklassen. Deswegen musst du auch in dieser Zelle bleiben.“

Da war Kjall aber nicht mehr einverstanden: „Bei Nehals verfaulten Zähnen, niemals bleibe ich hier zurück! Ich fliehe mit dir!“

„Das glaube ich nicht“, sprach der zweite Kjall und zog die schwere Zellentür zu, sodass Kjall in der Zelle feststeckte, während sein Doppelgänger außerhalb der Zelle stand. Ein leises Klonk verriet Kjall, dass sein Doppelgänger die Zellentür mit dem Casamatuc wieder verschlossen hatte. Kjall fluchte.

„Kjall! Lass mich hier nicht im Stich!“

Sein Doppelgänger entfernte sich ungerührt von der Zellentür.

„Wenn du mich nicht sofort hier herauslässt, schreie ich und alarmiere die Wachen!“

Der zweite Kjall drehte sich zurück: „Jetzt sei schon still. Wenn du dich geschickt anstellst, bist du schon bald in derselben Lage wie ich.“

Mit diesen Worten huschte der zweite Kjall in den dunklen, modrigen Seitengang zurück, in welchem die beiden Kjalls Bort und Hallwort zurückgelassen hatten. Kjall hörte seine Ketten ein letztes Mal klirren, dann war alles still.

Kjall trat lautstark gegen die Zellentür, einmal, zweimal, dreimal. Sie rührte sich keinen Millimeter. Zwergische Facharbeit. Es war zum Heulen. So nah an seiner Freiheit war er gewesen, und er hatte sich von einem Doppelgänger seiner selbst austricksen lassen! Nicht einmal sich selbst konnte man noch vertrauen. Konnte er sich vielleicht irgendwelche Vorteile davon verschaffen, das Versteck seines Doppelgängers zu verraten, sobald der Hallwort aus dieser Zeitlinie zurückkehrte? Und was hatte der zweite Kjall nur damit gemeint, dass Kjall selbst bald in seiner Situation sein könnte?


\newpage
\section{Geflohen}



Hektische Schritte kündigten die Rückkehr des Zwergenfürsten aus Kjalls Zeitlinie an. Kjall richtete sich erwartungsvoll auf. Hallworts breites Grinsen war durch das Gitterfenster in der Zellentür schon von weitem zu erkennen. Es verriet Kjall, dass der Fürst erfolgreich gewesen war. Triumphierend rannte dieser auf Kjalls Zelle zu und gab dann preis, was er in seinen Händen hielt: Eine etwa faustgroße Kugel aus umschlungenen Streben und kleinen Schaltern. Die Sphäre!

„Ich nehme an, das zeigt, dass meine Geschichte der Wahrheit entsprach“, setzte Kjall an, doch Hallwort unterbrach ihn:

„Mein lieber Jork, so sehr ich Eurem Wort auch Glauben schenken will, so fällt es mir nur schwer zu glauben, dass dieses abstruse Artefakt – so hübsch es auch aussieht – einen in eine andere Zeit bugsieren kann. Eine kleine Demonstration wäre bestimmt nicht zu viel verlangt. Welche Schalter muss ich betätigen?“

Kjall horchte auf. ‚Wenn du dich geschickt anstellst, bist du schon bald in derselben Lage wie ich...‘ Das hier könnte es gewesen sein, was der zweite Kjall damit gemeint hatte! Jetzt war Fingerspitzengefühl gefragt.

Kjall beugte sich eifrig vor und zeigte auf einige Querverstrebungen: „Wisst Ihr, dieses Gerät vermag es, ein Portal in Vergangenes zu öffnen. Dafür müsst ihr bloß die entsprechenden Querverstrebungen hier runterschieben, damit die Kontrollzelle in der Mitte freigegeben wird. Dort könnt ihr einstellen, wie weit in Vergangenes ihr reisen wollt... aber... aber nein, das ist eine Kunst, die sich nicht so mir nichts, dir nichts erlernen lässt. Nein, mein Fürst, haltet ein!“

Doch es war schon zu spät. Hallwort hatte in seiner Begeisterung den Haupthebel betätigt und die Maschinerie in Gang gesetzt. Die Sphäre begann, leise zu sirren. Blauer Dampf stieß aus einer Öffnung in ihrer Seite, waberte in die Luft und verformte sich zu einem immer schneller drehenden Portal nach Wer-weiß-wann. Hallwort ließ die Kugel erschrocken fallen und seine beiden Wachen stürzten nach vorne, um sich zwischen das Zeitportal und ihren Fürsten zu stellen. Sie sorgten sich umsonst, denn das Portal blieb nur wenige Augenblicke lang stabil, ehe es sich wieder zu einem kleinen, blau leuchtenden Funken zusammenzog und mit einem leisen „Puff“ verschwand.

„Mit Verlauf, mein Fürst, das war nicht geschickt!“, knurrte Kjall zwischen zusammengekniffenen Zähnen hervor. Wie viel Treibstoff befand sich wohl noch in der Sphäre? Falls die seltene Mischung aus flüssigem Feuer und Tiefsand vom Grund des Geheimen Sees, die der Sphäre ihre Kraft verlieh, wegen dieses neugierigen Trottels vollends verbraucht würde, würde Kjall ein für alle Mal hier feststecken. Und dabei war er so nah an seiner Flucht!

„Mein Fürst, ich würde vorschlagen, dass Ihr mich die Sphäre bedienen lässt“, schlug Kjall vor, „Ich könnte ein Portal wenige Minuten in die Vergangenheit öffnen, da kann bestimmt nichts schiefgehen.“

Hallwort blickte immer noch mit weit aufgerissenen Augen auf die Stelle, wo das Zeitportal verschwunden war. Dann fasste er sich und hob die Sphäre hoch, um sie an Kjall zu überreichen, als der Wächter Bort vortrat: „Nicht so schnell! Ihr werdet bestimmt nicht einfach so durch ein magisches Portal treten, welches von einem vorlauten Gefangenen geöffnet wurde.“

„Mein Fürst!“, flehte Kjall Hallwort an, „Ihr werdet durch das Portal hindurchsehen können und überprüfen, dass sich auf der anderen Seite nichts Gefährliches befindet. Spiegel in Vergangenes gibt es genug – habt Ihr nicht selbst einen solchen Hadrischen Spiegel in Euren Gemächern? Nur das Reisen durch das Portal wird beweisen, dass diese Kugel das ist, was ich behaupte. Und wolltet Ihr nicht schon immer die grundlegenden Gefüge der Welt erleben? Dies ist Eure Gelegenheit!“

Jetzt hab‘ ich dich, dachte Kjall, als Hallworts Augen zu glänzen begannen. Der alte Kindskopf. Stets auf der Suche nach einem Abenteuer.

Die zweite Zwergenwache trat nun ebenfalls hervor: „Ich muss doch tunlichst davon abraten. Und egal, was geschieht, ich verlange, dass der Gefangene während der Prozedur angekettet ist. Bort ist doch immer so stolz auf diese speziellen Eisenketten, die sein Vetter ihm geschenkt hatte. Das wäre endlich eine Gelegenheit, sie einzusetzen!“

„Genauso ist es“, setzte Bort gegenüber Kjall nach, „Ich halte mich nahe bei dir und lasse dich nicht aus den Augen, Freundchen. Und egal, was geschieht, diese Sphäre bleibt hier in diesem Raum in dieser Zeit! Nicht, dass du noch versuchst, mit ihr zu entkommen. “

Bort war nicht auf den Kopf gefallen, aber Hallworts Neugierde und Borts Respekt vor den Befehlen seines Fürsten würden dem entgegenwirken können. Dann blieb die Sphäre halt hier.

„Perfekt“, rief Kjall aus, „Dann kettet Bort mich an sich, ich öffne das Portal einige Minuten in die Vergangenheit und wir treten hindurch. Die Sphäre bleibt hier und ist somit außer Reichweite, falls ich sie wirklich stehlen wollte. Worauf warten wir noch?“

Jetzt musste es einfach schnell gehen, damit keiner sich Gedanken darüber machte, wie sich diese Zeitlinie hier weiterentwickeln würde, sobald das Trio durch das Portal getreten war. Hallwort war sicher nicht bereit, sein Reich spurlos zu verlassen. Ein Glück, dass zumindest die zweite Zwergenwache zu verwirrt schien, um sich einen Reim auf Kjalls Plan zu machen.\bigskip



Das Ganze war schnell abgewickelt. Kjall stand umgedreht an der Zellenwand, während Bort seine Hände und Füße mit Ketten verband und schließlich über eine längere Kette an seinem Gürtel befestigte. Der Schlüssel zu diesen Ketten legte er auf dem Tisch, weit von Kjall entfernt. Hallwort platzierte indes die Sphäre auf einem kleinen Strohhaufen auf dem Steinboden und trat dann zurück. Die zweite Schildwache musterte das ganze Geschehen verwirrt, während sie sich am Kopf kratzte.

Dann war es soweit. Kjall hantierte am Kontrollzentrum der Sphäre herum und stellte sicher, dass das Trio bloß eine Viertelstunde in die Vergangenheit reisen würden. Hallwort wollte sicherlich Kjalls Gesicht am anderen Ende des Portals erkennen, um sicherzugehen, dass es sie nicht in die tiefste Urzeit sandte, so aufregend ein Aufenthalt in der guten alten Urzeit Caverns auch wäre. Während bei Zeitreisen über Jahrzehnte oder gar Jahrhunderte hinweg zweifelsohne große Unsicherheiten ins Spiel kamen, konnte der Ankunftszeitpunkt für solche kurzen Zeitreisen glücklicherweise relativ klar bestimmt werden. Kjall überprüfte seine Einstellungen ein letztes Mal, seine Hände zitternd vor Vorfreude, und betätigte dann den Haupthebel.

Ein lautes Zischen ertönte, als flüssiges Feuer und Tiefsand sich mischten und blauer Dampf aus der Sphäre ausgestoßen wurde, welcher sich rasch zu einem Zeitportal formte. Diese Mal blieb das Zeitportal stabil. Auf der anderen Seite konnte man verschwommen erkennen, wie ein verschwommener früherer Kjall aufstand, sich dem Portal näherte und interessiert hineinblickte. Nun ja, genau genommen nicht der frühere Kjall, sondern ein Kjall aus einer anderen Zeitlinie. Oder?

Kjall blickte Bort und Hallwort an. Bort schaute finster zurück und brummelte etwas von wegen, dass Kjall es bereuen würde, wenn dem Fürsten auch nur ein Haar gekrümmt würde, während Hallwort nur noch Augen für das Portal übrig hatte.

„Auf drei“, flüsterte der Fürst.

„Bereit“, antwortete Bort. Kjall tat es ihm nach und begann dann, so unauffällig wie möglich seinen Atem zu verlangsamen. Er musste unbedingt als erster auf der anderen Seite des Portals erwachen.

„Eins!“, rief Hallwort.

Bort tatschte mit seinem Eisernen Handschuh nach dem Portal, welches unter seiner Berührung strahlend blau aufleuchtete.

„Zwei!“, rief Hallwort.

Kjall holte tief Luft.

„Drei!“

Die drei Zwerge traten synchron nach vorne und verschwanden ins Portal hinein.\bigskip



Kjall, Bort und Hallwort schwebten durch den endlosen Limbo, der von blauem Leuchten erfüllt war. Schwerelos hingen sie im leeren Raum, und doch bewegte sich die Gruppe langsam, unendlich langsam auf die helle Öffnung zu, die das andere Ende des Portals darstellte. Hallwort hatte den Fehler gemacht, zu versuchen einzuatmen, und keuchte nun verzweifelt, während er sich auf die Brust schlug. Bort hatte eine bessere Kondition und war ganz und gar nicht glücklich damit, wie die Situation sich entwickelte. Seine Handschuhe griffen nach Kjall, während er mit hochrotem Kopf versuchte, Kjall bei der Gurgel zu packen. Kjall fühlte sich, als würde sein Kopf explodieren, und kniff sich weiterhin stur die Nase zu. Hatten sie wirklich erst einen Drittel des Weges zurückgelegt? Hatte seine letzte Zeitreise durch den Limbo auch so lange gedauert?

Hallwort zuckte ein letztes Mal und erschlaffte. Bort ließ von Kjall ab und versuchte, sich durch den schwerelosen Raum auf seinen Fürsten zuzubewegen. Er schrie verzweifelt auf, ein Geräusch, welches stark verzerrt an Kjalls Ohren drang. Ohren, die mit schwerem Druck belegt waren. Nun zuckte auch Bort auf und erschlaffte. Kjalls Blickfeld wurde an den Rändern immer dunkler. Dann hielt er es nicht mehr aus, ließ seine Nase los und folgte seinen Instinkt, welcher ihn anschrie, einzuatmen. Glühender Schmerz stach in seine Lunge. Das gleißend blaue Licht wurde dunkler und verschwand dann endgültig. Das letzte, was Kjall wahrnahm, war ein metallisches Klonk, als die Kette, die ihn an Bort fesselte, auf Stein knallte.\bigskip



Ein metallischer Klirren war auch das erste Geräusch, das an Kjalls Bewusstsein drang, als er wieder aufwachte. Jemand hantierte hektisch an den Ketten zwischen seinen Füßen herum! Kjall riss seine Augen auf und blickte in sein eigenes Gesicht. Zwischen einem verfilzten roten Bart und Haarschopf blitzen kleine Augen hervor, welche ihn erstaunt und etwas erschrocken ansahen. Ein starkes Gefühl von Déjà-vu erfasste Kjall.

Kurzerhand packte Kjall den zweiten Kjall am Handgelenk, zog sich geschmeidig in die Höhe und blickte sich in der Zelle um. Das Zeitportal hatte sich bereits wieder geschlossen. Wie erwartet lagen Fürst Hallwort und Wächter Bort noch ohnmächtig neben ihm, während ein zweiter Kjall vor ihm stand und ihn anstarrte. Die Zellentür stand sperrangelweit offen und ein Casamatuc lag am Boden. Soweit, so gut. Jetzt musste es nur rasch gehen.

„Schnell“, rief Kjall seinem Doppelgänger entgegen, „Hilf mir, die beiden anderen aus dieser Zelle zu schaffen. Keine Sorge, das war ihre erste Zeitreise. So schnell werden sie nicht wieder erwachen.“

Ohne auf eine Antwort zu warten, packte der Kjall Bort an den Schultern und begann, ihn in Richtung der immer noch offen stehenden Zellentür zu bugsieren. Was hatte dieser Zwerg nur gegessen, schweres Roteisen?! Ein Schleifgeräusch und einige unterdrückte Flüche verrieten Kjall, dass sein Doppelgänger sich dem ohnmächtigen Hallwort zugewandt hatte und diesen aus der Zelle herauszerrte.

Kjall zog den ohnmächtigen Bort in einen dunklen Seitengang. Eine ehemalige Zelle? Ein Lagerraum? Jedenfalls lagen einige Decken am Boden. Kjall ließ Bort auf eine dieser Decken sinken und überlegte. Er musste so schnell wie möglich seine Ketten lösen und fliehen. Falls irgendwie möglich, im Besitz der Sphäre – so nahe würde er ihr so bald wohl nicht wieder kommen. In einigen Minuten würden sowohl der Schlüssel für seine Ketten als auch die Sphäre beinahe unbewacht in seiner Zelle liegen, wenn diese Zeitlinie sich gleich entwickelte wie diejenige, aus der er stammte. Und damit sich diese Zeitlinie gleich entwickelte, musste Kjall es nur seinem ersten Doppelgänger nachtun und diese zweiten Doppelgänger zurück in die Zelle sperren.

Ein lautes Ächzen zeigte Kjall, dass sein Doppelgänger soeben den Seitengang erreicht hatte und Hallwort erleichtert zu Boden plumpsen ließ. Kjall wandte sich an seinen Doppelgänger und sagte: „Wir müssen noch mal schnell in die Zelle zurück.“

Ohne auf eine Antwort zu warten, bewegte er sich in Richtung der Kerkerzelle zurück. Die Kette, die seine Hände und Füße an Bort knüpfte, war so lang, dass er gerade noch so die Zellentür erreichen konnte, ehe seine Kette anschlug.

„Der Casamatuc!“, zischte er seinem Doppelgänger zu, „Rasch!“

Der zweite Kjall huschte an ihm vorbei in die Zelle zurück, griff nach dem dort liegenden Casamatuc und übergab das Werkzeug an Kjall: „Er wird die Ketten leider nicht lösen können, das habe ich bereits versucht.“

Kjall nickte: „Ich ebenfalls. Aber wir wollen doch keine Spuren zurücklassen. Deswegen musst du auch in dieser Zelle bleiben.“

Auf dem Gesicht des zweiten Kjalls zeigte sich zunächst Verwirrung, dann rümpfte er die Nase: „Bei Nehals verfaulten Zähnen, niemals bleibe ich hier zurück! Ich fliehe mit dir!“

„Das glaube ich nicht“, sprach Kjall und zog die schwere Zellentür zu, sodass der zweite Kjall in der Zelle feststeckte, während er selbst außerhalb der Zelle stand. Geschickt setzte Kjall den Casamatuc an die Zellentür und drehte ihn leicht. Ein leises Klonk verriet ihm, dass die Tür sich erfolgreich verschlossen hatte. Sehr schön, alles lief nach Plan! Kjall drehte sich ungerührt um und entfernte sich wieder von der Zellentür, während sein Doppelgänger fluchte

„Kjall! Lass mich hier nicht im Stich! Wenn du mich nicht sofort hier herauslässt, schreie ich und alarmiere die Wachen!“

Dieses Wargorgehirn! Verstand er denn nicht, dass diese Situation ihm nicht schaden würde? Kjall seufzte. Natürlich verstand er es nicht. Noch nicht. Gereizt gab er zurück: „Jetzt sei schon still. Wenn du dich geschickt anstellst, bist du schon bald in derselben Lage wie ich.“

Mit diesen Worten huschte Kjall in den dunklen, modrigen Seitengang zurück, in welchem die beiden Kjalls Bort und Hallwort zurückgelassen hatten. Drei Knalle hallten durch den Gang. Sein Doppelgänger hatte soeben wütend gegen die Zellentür getreten.\bigskip



Bort wachte als nächster auf und war kurz davor, sich auf Kjall zu stürzen. Als er aber das kleine Messer sah, das Kjall von Hallworts Gürtel entfernt hatte, um es dem schlafenden Fürsten an den Hals zu halten, verhielt Bort sich plötzlich ganz ruhig. Kjall setzte dennoch warnend einen Zeigefinger auf die Lippen. Nicht, dass Bort noch auf den Gedanken kam, den Bort und den Hallwort aus dieser Zeitlinie mit seinem Geschrei zu warnen. Aber groß fürchtete sich Kjall nicht davor. Der Fürst bedeutete Bort viel zu viel, als dass er sein Leben so leichtfertig aufs Spiel setzen würde. Das Kjall noch nie mit einem solchen Messer umgegangen war und erst recht noch nie jemandem einen Todesstoß versetzt hatte, musste er dem Wächter ja nicht auf die Nase binden. Ein Mörder war Kjall bei weitem nicht. Aber einschüchternd wirken konnte er.

Hallwort regte sich gerade noch rechtzeitig, damit Kjall auch ihm wortlos bewusst machen konnte, in welcher Lage er sich befand. Laute Schritte aus dem dunklen Gang kündigten die Ankunft des Hallworts aus dieser Zeitlinie an, welcher gerade seine Sphäre zum zweiten Kjall in der Zelle brachte und am Lagerraum vorbeihastete, ohne die Zeitreisenden zu bemerken.

Hallwort sah sich interessiert um: „Ist das hier die Vergangenheit? Sieht nicht ganz so aus, wie ich es erwartet hätte.“

„Keinen Mucks mehr“, zischte Kjall den beiden Zwergen zu, „Wir wollen ja nicht eure Doppelgänger aus dieser Zeitlinie alarmieren.“.

Bort warf ihm nur einen hasserfüllten Blick zu und sah dann zu Boden. Hallwort hingegen gab sich völlig ungerührt angesichts des Messers an seiner Kehle und antwortete leise: „Wie kommt Ihr denn darauf, dass wir sie alarmieren könnten, werter Jork? Falls das überhaupt Euer Name ist – diese prekäre Situation lässt darauf schließen, dass Ihr uns in so einigen Dingen nicht die Wahrheit erzählt habt.“

Kjall ignorierte den Fürsten zunächst. Nach einigen Momenten der Stille überwiegte sein Interesse aber und er fragte: „Hallwort, wenn ich es wollte, was würde mich schon daran hindern können, jetzt hier rauszurennen und die Aufmerksamkeit deiner Wachen auf mich zu ziehen?“

Hallwort grinste ihn an, als er weiterflüsterte: „Nun, es gibt nichts, was Euch daran hindern würde, das zu tun. Aber ich weiß mit absoluter Sicherheit, dass Ihr es nicht tun werdet. Denn wenn Ihr es getan hättet, dann wären der gute Bort und ich ja alarmiert worden und befänden uns jetzt nicht hier. Es ist wie bei den echten Prophezeiungen – keine Kraft dieser Welt zwingt einen dazu, sie in Erfüllung zu bringen, und dennoch werden sie sich früher oder später bewahrheiten.“

„Kaulquappenquatsch!“, brummelte Kjall, „Das hier ist eine ganz andere Zeitlinie als die, aus der wir kommen.“

„Ach?“, meint Hallwort, „Und sie entwickelt sich zufälligerweise genau gleich wie die unsere, trotz unserem Einwirken? So lauscht doch...“

Von der Zelle am Gangende erklang soeben die laute Stimme Kjalls: „Ihr trinkt Euren Drachentee mit immer mit zwei Löffeln Honig drin, auch wenn Ihr das bei jeder passenden Gelegenheit abstreitet. Ihr führt eine geheime Tölplingzucht im Brauneisenstein. Als Kind seid Ihr einmal in den Geheimen See gefallen, als ihr heimlich und unerlaubt Cavern zu erkunden versuchtet, und eine von Wasser triefende, grünliche Gestalt rettete euch.“

Bort nickte nun auch: „Ich glaube, das sind dieselben Worte. Genau dieselben.“

Kjall schüttelte den Kopf: „Nein, nein, das muss eine andere Zeitlinie sein, sonst hätte ich das Vergangene ja gar nicht erst durch meinen Angriff auf die Andori verändern können!“

Hallwort antwortete nachdenklich: „Wie kommt Ihr darauf, dass Ihr sie verändert habt? Habt Ihr erreichen können, was Ihr erreichen wolltet?“

„Noch nicht“, knurrte Kjall, „Aber wartet nur, bis ich hier raus bin.“

Natürlich hatte Kjall sich schon die Möglichkeit überlegt, dass die Zeitlinie vorgeschrieben und unveränderlich sein könnte. Aber da er in keinem der unzähligen Geschichtsbücher und Schriftrollen, die er im Laufe seines Lebens verschlungen hatte, vom Angriff eines Schattens wie diesem Bleichen oder Shan oder von einem Angriff der Tiefminen-Golems auf die Brandurs Lager gelesen hatte, war er sich so sicher gewesen, dass die Sphäre ein Portal in eine andere Zeitlinie öffnen musste. Wie hätte er denn ahnen sollen, dass ihr Vorhaben so spektakulär scheitern würde? Nein, es konnte nicht sein! Es durfte nicht sein! Es musste eine Möglichkeit geben, Cavern vor dem Verfall zu bewahren, Xolls Tod zu verhindern und Radans Verbannung ungeschehen zu machen.

Hallwort unterbrach Kjalls Gedanken: „Wisst Ihr, ich habe mir bereits ausführlich Gedanken zu Zeitreisen gemacht. Einst kam ein Wanderer aus dem Hohen Hadria in Cavern zu Besuch, um seine Ausbildung zum vollwertigen Zauberer abzuschließen. Wie war sein Name schon wieder... Kohlaff? Er berichtete mir von Hadrischen Stundengläsern, die die Zeit verlangsamen konnten, und wir tratschten nächtelang darüber, was es denn bedeuten würde, wenn man die Zeit nicht nur komplett anhalten, sondern vielleicht sogar so sehr verlangsamen könnte, dass sie rückwärts abliefe. Wir kamen zum eindeutigen Schluss, dass Zeitportale, wie sie sich beispielsweise in den Tiefen des Horun zu manifestieren scheinen, einen nur in dieselbe Zeitlinie befördern können, aus der man stammt.“

Hallworts Augen glänzten schon wieder von Aufregung, und er schien das Messer an seiner Kehle komplett vergessen zu haben. Er strotzte vor Selbstbewusstsein, dabei musste er doch falsch liegen! Wenn Hallwort damit recht hatte, dass sie sich in derselben Zeitlinie befanden, aus der sie gekommen waren, dann könnten Kjall, der Fürst oder Bort ja mit jeder unbedachten Handlung Ereignisse in den Gang setzen, deren Auswirkungen früher oder später in einem Paradoxon enden würden! Was, wenn Kjall beispielsweise Xolls Tod verhinderte? Dann wäre Kjall selbst ja ganz anders aufgewachsen und hätte sich wohl nie dem Zeitreisen verschrieben – und der Kjall, der er in diesem Moment war, würde sich in Luft auflösen müssen, aufhören zu existieren!

„Und überhaupt“, fuhr Hallwort gedankenverloren fort, „Wenn Ihr mit diesem Portal in andere Zeitlinien reisen könntest, so würde das doch heißen, dass Ihr bei jeder Reise eine Zeitlinie zurücklässt. Wisst Ihr überhaupt, wie viele Zeitlinien dadurch bereits entstanden wären? Das Schicksal einer einzelnen Zeitlinie wäre doch komplett unbedeutend im See all dieser Möglichkeiten. Um das Beispiel vom Untergang Caverns zu nehmen: Ihr könntet vielleicht in eine Zeitlinie reisen, in welcher Cavern vor dem Untergang bewahrt wird, aber würdet dabei immer eine Zeitlinie zurücklassen, in welcher Cavern untergeht. Ich ging davon aus, dass mein zukünftiges Ich Euch solche Sachen erklärt hätte und dass der Untergang Caverns, aufgrund dessen ich Euch laut Eurer Behauptung hierhin zurück geschickt hätte, in der Zukunft bereits verhindert worden wäre...“

„Aber wenn es bereits verhindert worden wäre, warum würdet Ihr mich dann noch zurückschicken?!“, brauste Kjall auf. Sein Messer zitterte gefährlich nahe an Hallworts Adamsapfel, welcher erst jetzt wieder zu realisieren schien, dass er in Lebensgefahr schwebte.

„Nun, Jork, das müsst Ihr jetzt auch nicht sofort verstehen...“, versuchte Hallwort, Kjall zu beruhigen. Aber Kjall hatte genug. Er konnte seine angespannte Stimmung nicht mehr unterdrücken.

So leise wie möglich zischte er aufgebracht: „Weißt du, wie wir dich in der Zukunft nennen, wenn du erst einmal den Hammer abgegeben hast?! Hallwort den Nutzlosen! Hallwort den Pilzzüchter! Hallwort den Wahnwitzigen! Du hast nichts Großes an dir, du neugieriger Narr! Hältst dich für etwas Besseres als wir anderen. Lässt Cavern im Stich! Du steuerst unseren Untergang herbei, und du verstehst nicht einmal, dass es sich abwenden ließe! Ich sollte mich deiner hier und jetzt entledigen, immerhin würde Hallgard als Fürst nicht...“

„Das solltest du nicht tun!“, meldete sich Bort nun zum ersten Mal seit einigen Minuten wieder zu Wort, „Wenn du Fürst Hallwort etwas antust, wirst du diesen Lagerraum nicht mehr lebend verlassen, das schwöre ich bei der Axt meiner Urahnen!“

„Stimmt, das ist sehr spannend“, warf Hallwort ein. Seine Stirn war gefurcht. Offenbar war er nicht glücklich über die Spitznamen, die man ihm in der Zukunft verleihen würde, verkniff sich aber irgendwelche bissigen Antworten. „Offenbar wurden wir bei unserer Reise durch die Zeit auch örtlich versetzt, von Jorks Zelle in diesen Lagerraum, vermutlich aufgrund der natürlichen Drehung der Weltenscheibe...“

„Wieder falsch!“, fauchte Kjall nun, „Wir landeten am exakt selben Ort. Mein Doppelgänger und ich brachten euch ohnmächtige Zwerge in diesen Raum, ehe er in seine Zelle zurückkehrte!“

„Umso besser!“, rief Hallwort auf, „Das sollte meine Theorie doch bestätigen können. Habt Ihr bereits beide Rollen dieser Interaktion identisch durchlebt? Das wäre ein deutlicher Indikator, dass die Zeitlinie vorbestimmt...“

„Hallwort, wenn du nicht sofort den Mund hältst...“, setzte Kjall an.

„Drohe meinem Fürst noch ein einziges Mal, Verräter, und ich...“, fiel ihm Bort ins Wort.

„Vielleicht könnten wir alle einmal tief durchatmen“, schlug Hallwort vor, „Überlegen wir lieber, wie es weitergehen soll. Ich maße mir nicht an, Eure Motive zu verstehen, aber ich nehme an, dass Ihr gerne von hier fliehen würdet. Ich und Bort würden gerne weiterleben und könnten uns wahrscheinlich mit Eurer Flucht abfinden, solange Ihr fortan einen weiten Bogen um Cavern macht.“

Bort murmelte irgendetwas, das eindeutig nicht nach einer Zustimmung klang.

Kjall sammelte sich. Ja, er musste an seine Zukunft denken. Erst einmal weg von hier, raus aus Cavern, und dann könnte er noch lange genug überlegen, wie er den Verfall Caverns verhindern könnte.

„Das folgende wird nun geschehen“, setzte er an, „Wir bewegen uns alle drei rüber zur Zelle, wo inzwischen nur noch dieser zweite Schildwächter... wie lautet sein Name?“

„Ihr Name ist Ragga.“

„...wo sich inzwischen nur noch Ragga aufhält. Ihr entfernt meine Ketten, übergebt mir die Sphäre und ich ziehe von dannen. Falls ich irgendeine unbedachte Bewegung von Bort oder Ragga – oder Hallwort – sehe, so findet diese Messerklinge ein Ziel und Hallwort ein unschönes Ende. Wir wollen das alle nicht, also seht zu, dass es nicht soweit kommt. Alle einverstanden?“

Hallwort nickte ergeben, während Bort immer noch mit den Zähnen knirschte. Dann nickte er ebenfalls. Kjall empfand wieder Hoffnung. Bald hätte er das alles hinter sich.\bigskip



Ragga stand alleine inmitten der Zelle, hielt die Sphäre in ihrer Hand und versuchte wahrscheinlich gerade herauszufinden, wohin Hallwort, Bort und der Gefangene durch das Zeitportal verschwunden waren.

„Ragga, sei so lieb und roll die Sphäre zu mir rüber, willst du?“, fragte Kjall höflich, als sich seine Prozession der Zellentür näherte. Ragga blickte auf und erstarrte, als sie Kjalls Messerklinge an Hallworts Hals sah.

„Den Schlüssel zu meinen Ketten bitte ebenfalls. Ich will Bort nicht für den Rest meines Lebens hinter mir herschleppen.“

Eines musste man Ragga lassen, sie war erheblich schlauer als Bort und folgte Kjalls Anweisungen ohne Protest. Bort hingegen fluchte so sehr, dass Kjall sich gleich mehrere kreative Ausdrücke für spätere Gelegenheiten merkte, während Borts Ketten von Kjall gelöst und Bort damit an Kjalls Pritsche gekettet wurde.

Ein wenig Gerangel gab es tatsächlich auch noch, als Kjall die drei anderen Zwerge in der Zelle einzuschließen versuchte. Sein Messer an Hallworts Kehle war das einzige Druckmittel, das er hatte, und auch wenn das Leben des Fürsten der Schildzwerge einiges wert war, war sich Kjall natürlich bewusst, dass der kräftige Fürst ihn mit wenigen Faustschlägen zu Boden bringen konnte, wenn er ihm die Gelegenheit dazu gab. Also achtete Kjall ganz besonders auf Hallwort, als er diesen langsam losließ und ins Innere der Zelle dirigierte. Darum hätte Kjall beinahe Ragga übersehen, als diese sich überraschend auf ihn warf und ihn zu Boden drängen versuchte. Nur Kjalls schnelle Reflexe bewahrten ihn davor, für den Rest seines Lebens in einem Kerker der Schildzwerge zu versauern. So trat Kjall gerade noch rechtzeitig zurück, stieß Ragga von sich und riss – bereits zum zweiten Mal am heutigen Tag – die Zellentür zwischen sich und den anderen Zwergen zu. Mit der passenden Drehung des Casamatucs und einem dumpfen Klonk verschloss sich die Zellentür. Kjall trat zurück, Casamatuc in der linken, Sphäre in der rechten Hand.

Hallwort sah ihn nachdenklich an, Bort musterte ihn hasserfüllt und Ragga zeigte ihm eine unschöne Geste. Kjall hob seine Hand zum Gruß, drehte sich dann um und rannte den dunklen Gang entlang. Wahrscheinlich würde es einige Stunden dauern, bis die nächsten Wachen auf ihrer halbtäglichen Runde an dieser Zelle vorbeikommen und die gefangenen Zwerge bemerken würden, aber er wollte sein Glück nicht unnötig herausfordern. Erst recht nicht, als die Zwerge in der Zelle hinter ihm um Hilfe zu schreien begannen.

In der Stoffkleidung, mit denen die Schildzwerge ihren Gefangenen ausstatteten, schepperte Kjall nicht einmal, als er so schnell er konnte durch die finsteren Gänge und Stollen Caverns huschte, stets auf der Hut vor Patrouillen der Schildzwerge oder Minenarbeitern aus den Tiefminen. Obwohl sein Herz raste und er in jedem zweiten Schatten einen Zwerg erkannte, war Kjall in Hochstimmung. Er war nicht nur frei, er besaß sogar die Sphäre. Ihr kaltes Metall brannte in seiner Faust, und ein leises Summen verriet ihm, dass sie noch genug Saft für mindestens einen Zeitsprung besaß, wahrscheinlich sogar mehr. Nur wohin sollte er sich nun wenden? Die Stollen, die zu den Zellen führten, waren ihm unbekannt gewesen, aber diese Fackel von Cavern in der Halterung da vorne kannte er nur allzu gut! Kjalls Werkstatt und sein geheimes Museum lagen hier ganz in der Nähe!

Kurz überlegte Kjall, ob er seiner Werkstatt einen Besuch abstatten sollte, entschied sich dann aber dagegen. Wer weiß, wie sein jüngeres Ich darauf reagieren würde, wenn es einer faltigeren Version seiner selbst im Gefangenenoutfit gegenüberstehen würde? Lieber würde er zum Verlassenen Turm aufbrechen. Kjall wusste nicht, wie es seinen Verbündeten ergangen war, aber wenn der Bleiche, Shan oder Thogger die Begegnung mit den Helden besser überstanden hatten als er, würden sie sich wahrscheinlich ebenfalls dort einfinden. Die Schar hatte sich über längere Zeit in diesem Turm eingerichtet. Nur eine der unzähligen tragisch verlassenen Zwergenbauten im Grauen Gebirge. Das Zwergenreich war am Verenden und viele Zwerge kümmerte es nicht einmal! Kjall spuckte aus.\bigskip



Der Weg zum Verlassenen Turm verlief größtenteils ereignislos. Einmal musste Kjall einer Horde kichernder Zwergenkinder ausweichen, welche sich der Brücke zum Roteisenstein entlanghangelten. Wo waren nur deren Eltern?! Ein anderes Mal konnte er sich gerade noch rechtzeitig in einen Erker verdrücken, als eine Horde grölender Minenarbeiter an ihm vorbeizog. Und dieses Mal horchte er auf, denn einige dieser Leute erkannte er!

Da war der breite Radan, derjenige Zwerg, der für ihn wohl einem Freund am nächsten kam, eng versunken in eine energische Diskussion mit einem Spitzhackenträger. Auf dem Weg zu seinem Wachdienst? Noch hatte er keine Ahnung davon, dass er einst dafür verbannt werden würde, für seine Prinzipien einzustehen. Es war wirklich eine Schande.

Dort war der blonde Drak, der soeben prahlte vom köstlichen Waldpilz-Eintopf, den seine Gemahlin heute Abend für seine Familie zubereiten würde. Kram, einer von Draks unzähligen Kindern, war im Moment vielleicht sogar noch zu jung, um feste Nahrung zu kosten, sollte aber nach Hallworts Sohn Hallgard zum nächsten Fürsten Caverns werden und Radans Verbannung veranlassen... und... und steckte in diesem Moment ebenfalls in dieser Zeitlinie fest! Das hatte Kjall völlig vergessen! Der Kram aus der Zukunft befand sich ja mit den restlichen Helden auch hier in der Vergangenheit! Solange Kram ebenfalls in dieser Zeit festsaß, kümmerte es Kjall nicht einmal so sehr, falls er nicht aus dieser Zeit fliehen können würde. Ohne Kram, der noch keine Nachkommen, ja nicht einmal einen Lebensgefährten hatte, musste das Amt des Fürsten von Cavern doch an jemand anderes aus Hallworts Blutlinie fallen, der die Schildkrone zumindest ein bisschen mehr verdiente als ein dahergelaufener Zwerg aus den Tiefminen. Kjall musste nur darauf achten, dass die Helden mit ihrer Sphäre nicht seinen Plänen, Andor zu verhindern, einen Strich durch die Rechnung machen konnten.

Aber jeder Gedanke an Kram war vergessen, als Kjall das letzte Mitglied der Zwergenhorde erkannte: Da war Xoll Hammeraxt, Meisterschmied und Meisterkrieger – und außerdem Kjalls Vater. Kjalls Augen wurden feucht, als er sah, wie Xoll Hammeraxt den Stollen entlangstapfte und fröhlich über einen Witz Radans lachte. Erinnerungen an ihre gemeinsame Zeit schossen Kjall durch den Kopf. Konnte Kjall ihn davon abhalten, zur Befreiung der Rietburg aus den Händen dieses vermaledeiten Dunklen Magiers auszuziehen? Konnte Kjall Xolls unnötigen Tod verhindern? Was, wenn Hallwort recht gehabt hatte und Kjall immer noch in derselben Zeitlinie feststeckte – würde Kjall dadurch dann ein Paradoxon auslösen und seine eigene Existenz auslöschen? Nein, er hätte durch seine schiere Anwesenheit in dieser Zeitlinie doch bestimmt bereits mit irgendwelchen unbedachten kleinen Taten ein Paradoxon ausgelöst. Dass er noch hier war, bewies doch gerade, dass er, Hallwort und Bort sich in einer alternativen Zeitlinie befinden mussten. Und das hieß, dass er Xoll vielleicht vor seinem Tod bewahren können würde.

Diese und ähnliche Gedanken schossen durch Kjalls Kopf, während er durch die Gänge Caverns schlich, langsam, aber stetig dem Verlassenen Turm entgegen.\bigskip



Als Kjall an der großen Waffenkammer der Schildzwerge vorbeikam, hielt er inne und spähte vorsichtig hinein. Leer. Perfekt. Er schlich langsam im Schatten der Wände und Sockel durch den gewaltigen Raum, den die Zwerge einst ins Gebirge geschlagen hatten, und spähte umher. Feingearbeitete Kettenhemde, goldene Brustpanzer, große kantige Helme, alles fürstliche Rüstwaren, alles langweilig. Da fiel Kjalls Blick auf einen Schild in der Ecke. Er lief zu ihm, nahm ihn auf, betrachtete die reichen Verzierungen und dachte an Nehal, den Drachen und Kreatok, den Meisterschmied der Schildzwerge, die einst gemeinsam der Schmiedekunst nachgingen und die vier mächtigen Schilde erschufen. Dies hier war der Silberschild. Kjall lachte. „Dieses nutzlose Ding! Einer der vier ‚mächtigen‘ Schilde aus uralter Zeit, einfach so ungeschützt in einer Waffenkammer verstaubend! Ha, ha.“

Kjalls Sammeldrang überwog seine Furcht, durch den Diebstahl des Silberschilds ein Paradoxon auszulösen und seine eigene Existenz zu beenden. Diese Furcht legte sich endgültig, nachdem er den Schild angehoben und angelegt hatte, ohne sich in Luft aufzulösen. Als er den Verlassenen Turm erreichte, hatte er den Silberschild fest auf seinen Rücken geschnallt.\bigskip



Der Verlassene Turm war, nun ja, verlassen. Von Kjalls Kumpanen befand sich keiner dort, weder war die düstere Präsenz vom Bleichen oder Shan zu spüren, noch ließ sich der massige Tarus Thogger blicken. Kjall stand auf der Spitze des Turms und ließ seinen Blick über das Land schweifen, während die Sonne unterging und das Rietland in goldenen Farbtönen erglühen ließ. Dort drüben konnte er gerade noch den hohen Turm von Brandurs Lager erkennen. Nirgendwo war ein ungewöhnlich dunkler Schatten zu sehen. Waren seine Verbündeten gescheitert oder hatten sie sich nur an einen anderen Ort zurückgezogen? Es sah so aus, als wäre Kjall erstmals auf sich alleine gestellt. Aber das war in Ordnung. Die Schatten ließen ihm ohnehin immer Schauer über den Rücken laufen.

Was sollte Kjall als nächstes tun? Brandur entführen konnte er ohne seine Golems kaum, erst recht nicht, wenn die Helden aus der Zukunft ihn noch beschützten. Und weiter zurück in die Vergangenheit zu reisen, würde ihm auch kaum etwas bringen können, da er auf sich alleine gestellt nicht die Kampfkraft zur offenen Konfrontation besaß. Kjall wollte zurück in seine eigene Zeit, zu seiner Werkstatt, seinem Museum, zu seinen Tiefminen-Golems. Konnte er vielleicht seinem vergangenen Ich die Tiefminen-Golems stehlen, ehe dieses sie in einen zum Scheitern verdammten Angriff auf die proviorische Rietburg sandte? Wenn er diese erst einmal wieder hätte, könnte er noch tiefer ins Vergangene zurückreisen und Brandurs Schar vielleicht gleich nach ihrer Überquerung des Grauen Gebirges abpassen. Ironischerweise würde es Kjall sehr schwer fallen, die Kontrolle über die Tiefminen-Golems von seinem früheren Ich zu übernehmen, so lange die Tiefminen-Golems sein früheres Ich beschützen. Es war zum Verrücktwerden!

Kjall wollte im Moment nichts weiter, als sich auszuruhen, zu sammeln und an einem sicheren Ort einen neuen Plan auszuhecken. Aber einfach so zurückkehren in seine eigene Zeit konnte er auch nicht, da er nur wusste, wie er Portale in Vergangenes, aber nicht in die Zukunft öffnen konnte. Wobei... wenn er irgendwie wieder zurück in die Zukunft gelangen könnte, könnte sein zukünftiges Ich ein Portal in die Vergangenheit öffnen, an diesen Ort, zu diesem Zeitpunkt! Hoffnungsvoll blickte Kjall um sich, in der Erwartung, einen blauen Funken zu sehen, der sich zu einem Portal verformen würde. Kein Funke kam. Keine zukünftige Version seiner selbst öffnete ein Portal in die Vergangenheit, um Kjall vom Verlassenen Turm abzuholen. Was hatte das zu bedeuten?

Nach einer halben Stunde war die Sonne komplett untergegangen. Der blaue und der rote Mond standen beide hoch am Himmel, aber ihr Licht konnte gegen die Kälte der Nacht nicht ankommen. Aus der Ferne ertönte der kehlige Ruf eines Skrals. Kjall gab es auf, auf ein Portal aus der Zukunft zu warten, und zog sich ins Innere des Turms zurück, wo er mithilfe eines Glutstabs ein kleines Feuer entfachte. Dann wechselte er seine Gefängnismontur gegen eine verstaubte Wanderkleidung, die er nebst dem Silberschild und dem Glutstab in der Waffenkammer gefunden hatte. Der Silberschild! Kjall setzte sich neben den verzierten Schild mit den vier Stacheln und musterte ihn. Tastete ihn ab, versuchte, irgendeine magische Verbindung zu spüren. Nichts, außer dass der Nachtwind vielleicht noch etwas kühler durch das verlassene Gemäuer strich. Einige Zwerge nannten ihn den ‚Sturmschild‘ nach den mysteriösen Umständen um Hallworts Tod. Warum wollten diese Jungspunde für alles neue Namen finden? ‚Silberschild‘ war doch eine viel treffendere Bezeichnung.

Dieser Schild war für Kjall schon seit je her ein Rätsel gewesen. Was hatten sich Kreatok und Nehal nur dabei gedacht, als einen Schild ohne besondere Fähigkeiten erschaffen hatten? Eine Aussage über Ästhetik im Allgemeinen? Dass der Schild wunderschön war, ließ sich kaum bestreiten. Vielleicht waren die Fähigkeiten des Schildes einfach noch nicht entdeckt worden? Der Legende nach hatte er die Piraten mit Blitzen erschlagen, die ihn in einigen Jahren zu stehlen versuchen würden, aber ob das wirklich das Tun des Schildes gewesen war. Und ein letztes Mal reiterierte er: Kjall besaß jetzt den Silberschild. Demnach konnte Hallwort ihn in einigen Jahren nicht mit auf seine Reise in den Norden nehmen, und keine Piraten konnten versuchen, ihn zu stehlen. Kjall hatte durch den Diebstahl des Silberschilds nicht nur kein Paradoxon ausgelöst, sondern auch bewiesen, dass er mit seinen Handlungen die Zukunft ändern konnte und sich in einer neuen Zeitlinie befand. Oder? Hm... er musste sich morgen näher dazu Gedanken dazu machen. Jetzt wollte er einfach nur schlafen.

\newpage
\section{Gesehen}


Kjall konnte kaum schlafen. Er wälzte sich herum, auf den Bauch, auf den Rücken, auf die Seite, aber keine der Lagen war bequem. Nicht, dass er sich je groß um Bequemlichkeit geschert hätte. Er konnte sich jedoch nicht daran erinnern, derart heftige Schlafprobleme erlebt zu haben. Heute konnte er seine aufgewühlten Gedanken einfach nicht zur Ruhe bringen. Und so lag Kjall still da und ließ seine Gedanken schweifen.

Noch bevor die Sonne sich am nächsten Tag über die Berggipfel erhob, beschloss Kjall, diesem elenden Halbschlaf ein Ende zu bereiten und zum Baum der Lieder aufzubrechen. Wenn irgendwo weitere Schriftrollen zum Gebrauch dieser mystischen Sphäre zu finden waren, dann wohl in den Schwarzen Archiven. Ohne seine Tiefminen-Golems würde es für Kjall schwieriger werden, dort einzubrechen, aber das Risiko wollte er auf sich nehmen. Er wollte zurück in seine eigene Zeit, sein eigenes Museum, und wenn keine zweite Version seiner selbst ihn dorthin bringen wollte, musste er halt selbst herausfinden, ob diese Sphäre auch Portale in die Zukunft öffnen konnte. Schlimmstenfalls, falls er entdeckt würde, könnte er einfach ein kleines Stückchen weiter in die Vergangenheit fliehen.

Einige Apfelnuss-Sträucher am Rande des gewundenen Wegs verliehen Kjall etwas dringend gebrauchte Stärke, während er in seiner neuen Wanderkleidung durch den Wachsamen Wald strich, den Silberschild fest auf seinen Rücken gezurrt und mit einer dünnen Stoffdecke bedeckt, sodass er nicht allzu sehr auffiel. Kjall konnte nur hoffen, dass die Helden von Andor nicht ebenfalls beschlossen hatten, den Baum der Lieder aufzusuchen.\bigskip



Erstes Sonnenlicht ergoss sich über den Wachsamen Wald. Die Lichtung, über welcher der Baum der Lieder imposant thronte, lag vergleichsweise still da. Die meisten Bewahrer vom Baum der Lieder waren wahrscheinlich noch in ihren Betten oder kümmerten sich um ihre allmorgendlichen Rituale. Einige Dorfbewohner zogen bereits zum Brunnen und brachten Wasser zurück in ihre Häuser. Bogenschützen patrouillierten in Dreiergruppen am Rande der Waldlichtung entlang. Ein kleines Männchen saß auf einem Hocker am Rande der Lichtung und blickte Kjall aus zwei verschmitzten Augen unverwandt an. Wäre es vielleicht geschickter, umzudrehen und in der nächsten Nacht zurückkehren? Da fiel Kjalls Blick auf eine mächtige Gestalt, die in der Nähe des kleinen Männchens auf einem Stein saß, im Schneidersitz, die Hände auf die Knie gelegt und den Kopf mit den mächtigen Hörnern in Andacht gesenkt. Ein Tarus! Ein ihm wohlbekannter Tarus! Er schien etwas vor sich hin zu murmeln, und wenn Kjall sich nicht täuschte, wogte das Gras um den Stein herum, flimmerte die Luft um den Tarus leicht und knirschten die nahestehenden Bäume, als würde die Seele des Waldes dem Tarus antworten. Geisterhaft.

Um nicht den Verdacht der Wache haltenden Bogenschützen auf sich zu ziehen, näherte sich Kjall dem Tarus langsam, aber selbstbewusst aufgerichtet. Dass er abgesehen vom Schild auf seinem Rücken komplett unbewaffnet war, ließ ihn relativ ungefährlich erscheinen, dennoch wollte er keine Konfrontation auslösen. Das kleine bucklige Männchen betrachtete ihn weiterhin nachdenklich, während er näher schritt, nickte ihm dann aber freundlich zu und wandte sich von ihm ab. Er atmete erleichtert auf.

„Thogger“, zischte Kjall, und tippte dem Tarus auf die Schulter, die selbst im Sitzen die Schulter des Zwergs überragte.

Thogger schien nicht überrascht, sondern öffnete bloß die Augen: „Falls du dich anschleichen wolltest, so warst du eindeutig zu laut. Was willst du, Kjall?“

„Freust du dich nicht, mich zu sehen? Ich bin hier, um mehr vom Umgang mit der Sphäre zu erfahren. Ich will zurück in meine eigene Zeit reisen. Wie sieht es bei dir aus? Und wie steht es um die anderen?“

Thogger schnaubte: „Die haben versagt. Shan wurde in Brandurs Lager eingekerkert. Der Fahle hat tatsächlich einen Schleier der Nacht über das Land verhängt, aber dieser wurde kurz darauf wieder gebrochen, von einer Kräuterhexe an einem Feuerkreis. Ich wusste gar nicht, dass es einen in Andor gibt. Mit dieser Hexe und dem buckligen Hüter dort drüben habe ich mich schon ausführlich unterhalten. Sie verfügt über so viele Kenntnisse, dass sie eine gute Schamanin abgegeben hätte, und er ist weiser, als es selbst mein Vater Hogger war. Die Hexe scheint tatsächlich zu wissen, wo Sternkraut wächst, und könnte mich dorthin führen. Das ist gut, denn diese Helden, die mir dasselbe versprochen hatten, haben sich nach ihrem Sieg über den Fahlen in Luft aufgelöst. Ich vermute, dass sie einen Weg zurück in ihre eigene Zeit gefunden und uns hier zurückgelassen haben. Schöne Helden sind das. Ich weiß nicht, wie du ohne die Sphäre nach Hause zurückkehren willst.“

„Wenn es nur das ist...“, grinste Kjall, und ließ die schimmernde Sphäre aus einer seiner Taschen hervorblitzen. Thogger riss die Augen auf. Kjalls Grinsen wurde noch breiter: „Direkt aus der geheimen Fürstenkammer der Schildzwerge. Wenn du wüsstest, was ich alles erlebt habe... aber das ist jetzt nicht relevant. Kommst du mit, wenn ich herausfinde, wie wir uns zurück in unsere Zeitlinie transportieren können?“

Thogger zögerte: „Ich habe geschworen, mich Varatans Fluch hinzugeben, sobald das Sternkraut sicher in meiner Heimat angekommen ist. Diese Hexe vermag mich zum Kraut zu führen. Ich glaube, hier meine Bestimmung gefunden zu haben. Ich stärke mich nur noch für diese letzte Reise, dann breche ich auf. Meine Wünsche sind mit dir auf deinem weiteren Pfad, Kjall, aber ich glaube, hier trennen sich unsere Wege.“

Kjall nickte schwer. Er hatte nie ganz verstanden, was es mit diesen Flüchen auf sich hatte, die auf Thogger, Shan und dem Bleichen lagen, aber wenn Thogger hier seinen Frieden finden konnte, würde Kjall ihn nicht aufhalten. Interessanter war, dass die Helden offenbar einen Weg zurück in ihre eigene Zeit gefunden hatten. Eine Schande, dass Kram nicht in der Vergangenheit festsaß, aber diese Möglichkeit ließ in Kjall die Hoffnung wachsen, dass er dasselbe erreichen könnte. Zudem hieß das auch, dass die Helden ihm in dieser Zeit nicht mehr gefährlich werden konnten. Kjall musste nur den riesigen Baum betreten, die Treppe hochklettern, die Schwarzen Archive erreichen und die Schriftrollen finden, die ihm den Weg in die Zukunft zeigen würden. Kein Problem, er hatte bereits einem Fürsten und seinen zwei Wachen getrotzt, was sollten ihn da einige blasierte Schriftensammler auch aufhalten.

Kjall griff dem riesigen Tarus an den Unterarm und sprach ihm viel Glück und den Segen des feurigen Gottes für seine Zukunft zu. Der Tarus verabschiedete sich gerührt und Kjall bewegte sich schnurstracks über die Lichtung auf den Baum der Lieder zu. Vielleicht war es doch eine schlechte Idee gewesen, bei Tageslicht hierher zu kommen. Die Bogenschützen hatten ihn nicht aufgegriffen, vielleicht konnte er aufgrund seiner Waffenlosigkeit auch als harmloser Gelehrter durchgehen, aber sein Aussehen weckte dennoch die Aufmerksamkeit vieler Dorfbewohner. Zwerge sah man hier nicht alle Tage, und erst recht nicht welche, die sich mit dem Tarus unterhielten, der erst kürzlich noch Sturm und Unwetter über den Baum beschworen hatte. Und der Silberschild auf seinem Rücken sah selbst von einem Tuch verdeckt immer noch imposant aus.

Ich bin ein ganz unschuldiger, schriftrollenversessener Narr, ganz genau wie alle anderen hier, dachte Kjall, achtet einfach nicht auf mich, ich bin nur wegen der Schriftrollen hier. Er schaffte es nur bis kurz vor den Eingang zum großen Tor am Baum der Lieder, dann hörte er eine helle Stimme vom Rande der Lichtung: „Das ist Jork, dieser elende Zwerg, der in die Schwarzen Archive eingebrochen ist!“

Kjall sah sich unauffällig um – oder besser gesagt versuchte es zumindest – und erblickte eine junge Bewahrerin, die mit ausgestrecktem Finger auf ihn zeigte. Er hatte diese Frau noch nie in seinem Leben gesehen.

Thogger sah nun ebenfalls alarmiert auf und meinte beschwichtigend: „Aber was redest du denn da, Folla, das ist doch nur Kjall, ein alter Freund...“

Folla quittierte diese Aussage mit „Schweig schon stille, unseliger Unruhestifter!“ und beharrte weiterhin darauf: „Das ist Jork, der Mörder von Fanatos! Fasst ihn!“

Das brachte Leben in die drei Bogenschützen, die Folla am nächsten standen. Kjall blieb nicht lange genug stehen, um ihnen beim Anrennen zuzuschauen. Stattdessen nahm er selbst die Beine in die Hand und hechtete auf den Eingang zum Baum der Lieder zu. Links und rechts von ihm zischte je ein weiß befiederter Pfeil durch die Luft, und zwei dumpfe Geräusche an seinem Rücken verrieten ihm, dass zwei weitere Pfeile am Silberschild an seinem Rücken abgeprallt waren. So war dieser Schild zumindest für etwas gut. Dann hatte Kjall das Innere des Baums der Lieder erreicht.

Ein Glück, dass er den inneren Aufbau des Baums so gut studiert hatte, als er seinen ersten Einbruch mit den Tiefminen-Golems geplant hatte. Dieses Meisterwerk aus Architektur und Naturkunde, eine Vielzahl von Räumen und Wohnungen in einem lebendigen Baum, hatte Kjall stark an die Art und Weise erinnert, wie die Zwerge Gänge und Stollen durchs Graue Gebirge gruben. Auch hier fühlte er eine tiefe Verbundenheit mit der Natur, als er im Innern des riesigen Baums nach oben sah. Aber davon durfte er sich nicht ablenken lassen.

Kjall wusste genau, wie er die Schwarzen Archive erreichen konnte. Die Wendeltreppe mit den viel zu großen Stufen hoch, durch den Eingang auf die Balustrade hinaus, um die halbe Balustrade herum, die Raumflucht betreten, zweite Tür links. Auf seinem Weg traf er auf allerlei verschlafenere und wachere Menschen, welche ihm neugierig hinterhersahen. Novizen, Adepten, untere Bewahrer, selbst der Oberste Bewahrer... Kjall hatte für sie keine Zeit, denn sobald die Bogenschützen unter der Anführung Follas die Wendeltreppe erreicht hatten, zischten wieder Pfeile an ihm vorbei.

„Fasst ihn! Kallun, so schnapp‘ ihn dir doch! Er hat Fanatos auf dem Gewissen!“, ertönte erneut Follas Stimme, und die ersten Bewahrer versuchten, sich in Kjalls Weg zu stellen. Einer bekam ihn auch tatsächlich zu fassen. Als Kjall sich losriss, löste sich das Tuch von seinem Rücken und enthüllte den Silberschild in all seiner Pracht. Dabei kannte Kjall doch gar keinen Fanatos. Von was sprach diese Frau nur?

Endlich erreichte Kjall die Raumflucht, die die Schwarzen Archive enthielt. Gerade noch rechtzeitig schaffte er es, die Tür hinter sich ins Schloss fallen zu lassen, als von draußen auch schon die ersten Faustschläge zu hören waren.

„Gib auf, Jork!“, schrie Folla nun, „Nur wenn du dich freiwillig stellst, wird das Urteil des Rats der Bewahrer milder ausfallen! Und der Rat wird dich früher oder später richten, du Mörder, da kannst du dich darauf verlassen!“

Ein flaues Gefühl machte sich in Kjalls Magen breit, aber er achtete nicht darauf. So war das nun wirklich nicht geplant gewesen. Aber wie schon bei ihren Anschuldigungen, hatte diese Folla auch mit ihrer letzten Aussage hatte Unrecht. Kjall verfügte sehr wohl über die Möglichkeit, dieser Situation ungestraft zu entkommen.

Nachdem Kjall mithilfe des fürstlichen Casamatucs sicherstellte, dass die Tür zu den Schwarzen Archiven demnächst verschlossen blieb, beschloss er, auf Nummer sicher zu gehen, und wuchtete ein schweres Schreibpult vor die Tür. Viel Zeit würde ihm das nicht verschaffen, aber zumindest genug, um die passende Stelle im Schwarzen Archiv aufsuchen zu können.

„Zwergenartefakte, Zwergenartefakte...“, murmelte Kjall, während er die langen Regale auf, und abwanderte. Das Archiv war nach einer gewissen Logik sortiert, die sich ihm noch nicht zur Gänze erschloss. „Wo sind denn die verflammten Zwergenartefakte?!“

„Den Gang herunter, zweites Regal links“, ertönte eine heisere Stimme. Kjall schreckte auf. An einem der Schreibpulte zu seiner Linken saß eine uralte Frau, in ein schneeweißes Gewand gekleidet, welche mit einem langen Finger zitternd den Gang entlangzeigte. Kjall starrte die eigenartige Frau an. Sie war unverwechselbar eine Bewahrerin, und eine altehrwürdige dazu. Warum half sie ihm? Ihr Gesicht blieb abgewandt, als sie sanft sprach: „Nun zieht schon von hinnen, Störer des Friedens. Ihr werdet nicht zufrieden sein mit dem, was ihr vorfindet. Die Schriften sind nicht mehr hier.“

„Nicht mehr hier?“, fragte Kjall, „Wo sind sie dann?“

„Wenn wir das wüssten, mein Kind, wenn wir das wüssten. Sie wurden uns gestohlen.“

Dann drehte die Frau ihm endlich ihr Gesicht entgegen. Ihr beinahe zahnloser Mund verzog sich zu einem Lächeln und blinde Augen starrten ins Nichts, als sie prophetisch meinte: „Du wirst hier keinen Frieden bringen, Kjall, Sohn des Xoll, und du wirst keinen Frieden finden. Dein Ende wird langsam kommen und schmerzvoll sein, wie du es verdient haben wirst. Und doch... wenn du dich je dazu entschließen solltest, deine Geschichte zu Pergament zu bringen, so würden wir uns geehrt führen, sie hier zu bewahren.“

Kjall starrte die alte Frau entgeistert an und wich verängstigt zurück. Erst ein weiteres lautes Klopfen vom Eingangstor zu den Schwarzen Archiven holte ihn in die Gegenwart zurück. Jetzt hörte er sogar einen Schlüssel im Schloss knirschen.

„Alte Hexe“, fauchte er der seltsamen Frau entgegen, und rannte dann in Richtung des Ganges, auf den sie gezeigt hatte. Zweites Regal links, zweites Regal links... da war es!

„Möge Irlok dich holen!“, stieß Kjall hervor, als er endlich die Ablage fand, die in großen andorischen Buchstaben mit „Zwergenartefakte“ betitelt war. Hier hatte er gehofft, weitere Informationen zur Sphäre zu finden. Aber im Gegensatz zu sämtlichen sonstigen Ablagen in den überladenen Regalen dieses Archivs war diese Stelle vollkommen leer. Keine einzige Schriftrolle zu den Zwergenartefakten befand sich in diesem Regal. Die alte Bewahrerin hatte nicht gelogen. Aber das war jetzt nicht die Hauptsache. Wichtiger war, diesem verrückten Mob aus aufgebrachten Bewahrern zu entkommen.

Kjall packte seine Sphäre aus und legte sie vor sich auf den modrigen Archivboden. Wie weit sollte er in die Vergangenheit reisen? Eine Woche, einen Mond, gar ein ganzes Jahr? Falls er nicht herausfände, wie er in die Zukunft reisen könnte, würde ihn das alles nur tiefer in die Vergangenheit schicken. Und tiefer in der Vergangenheit... die Zeiten, in denen Cavern noch glorreich und voller Tradition gewesen war... ach, es wäre schon zu schön, diese Zeiten zu erleben. Aber es war wichtiger, sich um die Zukunft des Zwergenreichs zu kümmern, und das konnte er tief in der Vergangenheit nicht. Rund ein Mond zurück musste reichen. Das würde ihm nebenbei, falls er hier scheitern und nicht in die Zukunft reisen sollte, genug Zeit verschaffen, in Ruhe einen neuen Plan auszuhecken, ehe sein Doppelgänger aus der Zukunft mit Thogger, Shan, dem Bleichen und den Helden im Gepäck hier erscheinen würde.

Das Kontrollzentrum der Sphäre war rasch bedient und das Zeitportal entstand in Kürze. Kjall blickte hindurch und sah durch die blau schimmernde Scheibe – wenig überraschend – denselben Regalabschnitt, den er auch so hinter dem Portal sehen konnte. Personen waren keine zu sehen. In diesem Moment ertönte ein lautes Krachen von der Tür zu den Schwarzen Archiven und Follas Stimme: „Gut gemacht, Kallun!“

Die Stimme der alten Bewahrerin ertönte: „Wirklich toll gemacht, Kallun! Der Kerl ist doch schon so gut wie weg und neue Türen wachsen auch nicht auf Bäumen. Mhare, mein Kindchen, sei bitte so gut und begleite deine Großmutter zurück auf ihr Zimmer. Das war mehr als genug Aufregung für heute. Und ohnehin, Folla, dein Vorgehen...“

Kjall vernahm nicht mehr, was die alte Bewahrerin an Follas Vorgehen auszusetzen hatte, denn er hatte sich bereits die Sphäre gegriffen und war damit ins blaue Zeitportal gesprungen.\bigskip



Der Geruch nach altem Holz stieg in Kjalls Nase. Kjalls Nase, welche auf dem Archivboden flachgedrückt wurde. Der Archivboden des Schwarzen Archivs, auf dem Kjall ziemlich verrenkt ruhte und jeden Moment entdeckt werden konnte!

Rasch richtete Kjall sich auf und stieß prompt mit dem Kopf an eine Tischkante. Er unterdrückte einige Flüche und sah sich um. Aus einem vergitterten Fenster brach Sonnenlicht ins Archiv und tauchte den Raum in ein warmes, goldenes Licht. Staubfetzen tanzten durch die Luft und eine Fee zog sich laut meckernd vom Fensterrahmen zurück, Kjall habe mit seinem unordentlichen Magiedingsbums ihren Schönheitsschlaf gestört, ein Verbrechen, für das der Rat der Bewahrer bestimmt die Höchststrafe verhängen würde.

Kjall wusste besser, als sich mit einer Fee anzulegen, also brummelte er eine leise Entschuldigung und wandte sich ab. Sein Blick fiel auf die Ablage, die mit „Zwergenartefakte“ angeschrieben war. In der Zeitlinie, aus der er soeben gekommen war, war sie vollkommen leer gewesen. In dieser Zeitlinie, einen Mond früher, war sie allerdings nur so überfüllt mit Manuskripten, Schriftrollen, Pergamenten aller Art. Kjall hatte seine Tiefminen-Golems bereits zu einem höchst riskanten Einbruch in die Schwarzen Archive gesandt, aber diese Dokumente hier hatte er noch nicht studieren können. Diese Gelegenheit konnte er sich nicht entgehen lassen. Auf der Suche nach einem Gefäß drehte sich Kjall im Kreis. Eine alte Holzkiste lag da unter einem Pult, von Spornenweben übersät. Ihr Deckel fehlte und außer Staub schien sie nichts zu beinhalten. Das musste genügen.

Kjall wuchtete die Holzkiste unter dem Lesepult hervor und stand auf die Zehenspitzen, um die „Zwergenartefakte“-Ablage zu erreichen. Ein Glück, dass sie so tief unten stand! In Windeseile zerrte Kram wahllos Schriften und Berichte aus der Ablage in die verstaubte Holzkiste. Egal, was diese Berichte beinhielten, er konnte nur davon profitieren.

„Jetzt stört er nicht nur meine Ruhe, sondern stiehlt auch noch seltene Schriften. Ts, ts, was der Oberste Bewahrer nur davon halten würde?“, ertönte nun die neckische Stimme der nervigen Fee, die sich vom Fenstersims erhoben hatte und Kjalls Kopf umschwirrte. Sie fuchtelte mit ihren kleinen Händchen herum und Funken sprühten drohend aus ihren Fingern.

„Ich klau‘ doch nichts, ich bin nur hier, um zu studieren!“, erwiderte Kjall unwirsch, wedelte die glühenden Teilchen beiseite und kletterte demonstrativ auf den Stuhl vor einem Lesepult. Der Silberschild auf seinem Rücken drückte unangenehm gegen die Stuhllehne, aber das war jetzt nicht wichtig. Einer der Berichte, die Kjall von der Ablage in die Holzkiste geschaufelt hatte, hatte sein Interesse ganz besonders erweckt. Dieses grau-silberne Papier stammte eindeutig aus der Fürstenkammer Caverns... und er hatte bislang nur einen einzigen Bericht gelesen, welcher auf diesem Material geschrieben war.

Rasch kippte Kjall die restlichen Schriftrollen der „Zwergenartefakte“-Auslage in seine Holzkiste und wandte sich mit zitternden Fingern dem Bericht auf dem grau-silbernen Papier zu. Er traute seinen Augen kaum, als er eine weitere Skizze eines ihm so wohlbekannten Zwergenartefakts erkannte. Kein Zweifel, er war auf Gold gestoßen! Der zweite Teil eines Berichts zur geheimnisvollen Sphäre!

Wie schon der erste Teil des Berichts, den Kjall aber und abermals konsumierte hatte, war auch dieser hier in schlichten Lettern, kurz und bündig geschrieben. Wer weiß, ob hier vielleicht mehr Informationen zum Erschaffer der Sphäre zu finden waren. Im Moment musste Kjall nur eines: Herausfinden, ob und wie er zurück in die Zukunft reisen konnte.

„Meint der Zwerg etwa, er könne mich ignorieren?!“, erklang erneut die quengelnde Stimme der Fee, die sich auf Kjalls Lesepult niedergelassen hatte. Noch mehr Funken tanzten um die Fee herum und drohten, die nahegelegenen Schriftrollen zu entzünden.

Die Fee fuhr fort: „Ich weiß, ich weiß, ich soll‘ ja um Mutters Willen nicht die edlen Herrschaften stören, die sich hier um die Bewahrung der Geschichten und Legenden des Landes kümmern, das wurde mir oft genug eingebläut. Aber weiß er, ich hab‘ da so den gewissen Verdacht, dass er gar nicht zu den edlen Herrschaften gehört, die diesen Raum betreten sollten. Ich hab‘ einen Riecher für so etwas.“

Kram guckte die Fee enerviert an und erkannte zum ersten Mal, dass es sich um einen Feerich handelt. Der Feerich tippte sich auf die Nase und kicherte, als er Kjalls Gesichtsausdruck sah, und meinte: „Ertappt! Wie spannend, sonst geschieht hier fast nie etwas aufregendes! Was sucht er hier? Will er geheimes Wissen klauen? Wie gedenkt er, von hier zu entkommen? Soll ich die Wachen alarmieren?“

Kurz überlegte Kjall sich, einfach einen weiteren Mond in die Vergangenheit zu reisen, um diese Fee hinter sich zu lassen. Aber er wusste nicht, wie viel Treibstoff noch in der Sphäre lagerte, und er würde diesen sicher nicht wegen eines vorlauten Feerichs verbrauchen. Er musste nur diesen Bericht überfliegen... da! Da stand eindeutig etwas von einer Reise in ‚Zukünftiges‘. Und dafür musste man... oh, natürlich! Die Querverstrebung links und der dritte Schalter im Kontrollzentrum. Nun, jetzt war es offensichtlich! Warum war Kjall nicht von selbst darauf gekommen?

„Meint er etwa, er könne mich einfach ignorieren? Hallihallo, mich gibt es auch noch!“

Der Feerich flatterte vor Kjalls Gesicht und schnippte mit den Fingern. Der Bericht zur Sphäre, den Kjall ehrfürchtig in seinen Händen hielt, verging in einem Funkenregen. Kjall zog seine Hände abrupt zurück und betrachtete fassungslos die verkohlten Papierfetzen, die vor wenigen Augenblicken noch ein Bericht zur Sphäre gewesen waren. Dann wandte er seinen hasserfüllten Blick dem Feerich zu. Der Bericht zu einem der mächtigsten Artefakte aus aller Zeit, vernichtet durch die Willkür einer Fee! Funken und Feuer, was für ein nutzloses Wesen! Kjall sah rot.

„Ups!“, grinste der Feerich, „Habe ich nun seine Aufmerksamkeit? Die Wachen werden in wenigen Augenblicken hier oben sein. Ich glaube, damit ist seinen Schandtaten ein Ende...“

Kjall griff nach seinem Casamatuc und versetzte dem kleinen Feerich einen mächtigen Hieb. Der Feerich reagierte nicht rechtzeitig und wurde aus der Luft an das nebenstehende Regal gefegt, wo er gebrochen zu Boden sank. Kjall musste ihn nicht genauer ansehen, um zu wissen, dass er das nicht überlebt haben konnte.

‚Mörder‘, klang Follas Stimme aus seiner Erinnerung nach, ‚Du hast Fanatos auf dem Gewissen!‘

Schwer atmend stand Kjall von seinem Pult auf. Hatte Folla recht behalten? War er ein Mörder? Aber... aber es war doch nur eine elende Fee! Kjall hatte noch nie einem Zwerg oder selbst einem verflixten Menschen etwas angetan und vermochte dies wahrscheinlich auch gar nicht! Klar, Brandur hätte er mit Freuden entführt, aber das wäre für das größere Wohl Caverns gewesen. Und dieser Feerich, dieser Nervtöter, hatte nicht einmal ansatzweise eine Ahnung davon, welches Schriftstück er soeben vernichtet hatte. Nicht... nicht dass er deswegen verdient hätte, zu sterben.

Bei Kreatoks versengten Augenbrauen, Kjall wusste ja jetzt, wo sich der Bericht befand! Er würde einfach hierhin zurückreisen und den restlichen Bericht lesen können. Es war nicht von Bedeutung, weder die Vernichtung des Berichts noch den Tod dieses Störenfrieds hier. Es war einfach nur ärgerlich. Aber gut, nun musste Kjall definitiv von hier verschwinden. Zum Glück wusste er jetzt auch, wie. Die Querverstrebung links und der dritte Schalter im Kontrollzentrum. Ganz einfach.

Als wäre das ihr Stichwort gewesen, hörte Kjall, wie sich die Tür zu den Schwarzen Archiven mit einem lauten Knarren öffnete. Eine helle Stimme, die er als Follas erkannte, rief: „Fanatos, ist alles in Ordnung? Warum hast du uns gerufen?“

Die Besorgnis in Follas Stimme wurde zu Gereiztheit, als die Bewahrerin nachsetzte: „Ich schwöre, wenn das wieder einer deiner Scherze ist...“

Kjall brachte die Sphäre aus seiner Tasche zum Vorschein und verschob hastig Streben, Schalter und Hebel. Keine Zeit, den genauen Zeitpunkt zu bestimmen, einfach in die Zukunft zurück wollte er. Lieber etwas zu weit, wenn er bedachte, dass er sich soeben mit dem gesamten Bewahrerorden zu verfeinden schien. Kjall platzierte die Sphäre in der Holzkiste mit den gestohlenen Schriftrollen und drehte sie leicht, sodass das Portal sich seitlich um die Holzkiste bilden würde. Diese Kiste würde Kjall sofern möglich gerne in die Zukunft mitnehmen.

Die Sphäre begann zu sirren und blauer Dampf stob aus einer Öffnung zu ihrer Seite, während die drohenden Schritte Follas immer näher kamen. Zu wenig Zeit, zu wenig Zeit! Kurzerhand wandte sich Kjall vom sich bildenden Zeitportal ab und trat selbstbewusst zwischen den Regalen hervor, den Silberschild immer noch auf seinen Rücken geschnallt.

„Zum Gruße, werte Bewahrerin“, rief Kjall, „Ich bin Jork von den Schildzwergen“ – er tippte auf den Schild an seinem Rücken – „und ich würde hierher gesandt, um... um... um die Stabilität dieser Räume zu begutachten.“

Folla blieb überrascht stehen, nur wenige Schritte von Kjall, dem toten Feerich und dem sich bildenden Zeitportal entfernt. Es dauerte nur einige Augenblicke, bis sie zwei gekrümmte Klingen aus ihrem Gürtel gezogen hatte und Kjall entgegenstreckte. Immerhin hatte sie das Portal und Fanatos‘ Leichnam noch nicht entdeckt.

„Wie seid Ihr hier hereingekommen, Jork? Zu diesem Raum besitzen nur wenige Auserwählte Zugang...“

„Ich muss sagen, ich bin wirklich beeindruckt von den Abstützungen, die diese Deckenkonstruktion verstärken. Wer auch immer diese Räume in den Baum gebaut hat, die wussten wirklich, was sie taten“, plapperte Kjall munter weiter, immer wieder zum sich bildenden Zeitportal schielend. Inzwischen hatte es volle Größe erreicht und seine Rotation nahm stetig zu. Nicht mehr lange...

„Was wollt Ihr hier?“, fragte Folla nun, „Der Oberste Bewahrer würde bestimmt nicht einfach... nun, fragen wir ihn doch einfach selbst. Wenn Ihr mir bitte folgen würdet? Dem Geheimnis Eures Aufenthalts hier werden wir noch auf den Grund gehen, glaubt mir! Aber das müssen wir nicht in hier in diesen geheimen, zutrittsverbotenen Räumen tun.“ Den Worten „geheim“ und „zutrittsverboten“ verlieh sie eine besondere Betonung.

„Ich befürchte, dass ‚zutrittsverboten‘ kein wirkliches Wort ist. Und ich befürchte, dass ich euch nicht folgen werde“, antwortete Kjall vorsichtig. Aus dem Augenwinkel sah er, wie die Holzkiste mit den Schriftrollen und der Sphäre darin ins schief stehende Portal kippte. Der Durchgang war offen!

„Wenn ihr euch nicht gefügig verhaltet, so sehe ich mich gezwungen...“, begann Folla, doch Kjall achtete nicht mehr auf sie. Er drehte sich auf seinen Fersen um, rannte zwischen die Regale und stürzte auf das Portal zu. Ein scharfer Schmerz entflammte in seinem rechten Arm, dann hatte er das blau schimmernde Portal erreicht und ließ sich hineinfallen. Kalt glitt dessen Oberfläche über seine Haut, während die üblichen blauen Strahlen sein Blickfeld überdeckten. Das letzte, was Kjall wahrnahm, war ein gellender Schrei Follas.

„Fanatos!“\bigskip


\az{Jahr 340}


Kjall erwachte nießend in einem riesigen Berg aus Staub und zerfledderten Pergamenten. Kein Tageslicht und keine Fackeln erhellten das Innere des Raums. Nur schwaches, goldenes Mondlicht drang durch das Fenster zu seiner linken. Blind tastete Kjall seine Umgebung ab und ergriff die Holzkiste mit zahlreichen Schriften zu verschiedenen Zwergenartefakten. Er rang nach Luft. Warum war es hier so schwer zu atmen?! Natürlich war die Lage hier kein Vergleich zum schrecklichen Limbo zwischen der Zeit, aus dem er soeben gekommen war, aber dennoch kam Kjall mit der schlechten Luft nur schlecht klar. Endlich trafen seine Finger auf kühle, glatte, verschlungene Metallstreben! Kjall ergriff die Sphäre, steckte sie in seine Tasche und hechtete zum nahe gelegenen Fenster.

Kjall musste auf den Sims klettern, um das Fenster zu ergreifen. Er riss und zog, und unendlich langsam ruckelte das Glas auf. Kühle, frische Luft strömte in die Schwarzen Archive, und Kjall sog sie gierig ein. Gute Güte, wie lange hatte niemand diese Archive betreten? Das war falsch, das sollte so nicht sein. Kjall hatte es offenbar in eine weit entfernte Zukunft verschlagen, erheblich weiter, als er geplant hatte.

Ein frischer Schnitt brannte an Kjalls rechtem Arm. Offenbar hatte Folla ihn mit einem ihrer Dolche erwischt. Aber tief war er nicht und die Blutung hatte bereits gestoppt. Alles im grünen Bereich.

Zwischen den schneebedeckten Ästen des Baums der Lieder hindurch – es musste Winter sein, was auch die verdammte Kälte erklären würde – sah Kjall in der Ferne Lichter aufblitzen. Konnten das der Sternenhimmel über den nördlichen Ausläufern des Grauen Gebirges sein? Irrte er sich, oder waren weniger Sterne am Himmelszelt zu sehen als üblicherweise? Kjall schüttelte den Kopf und wandte sich wieder vom Fenster ab.

Am Lesepult fand er eine schon fast vollständig abgebrannte Wachskerze in einem Tonhalter, welche er mit einer beiläufigen Bewegung seines Glutholzes anzündete. Das Licht der Kerze enthüllte erst jetzt den Zustand der Schwarzen Archive. Viele Schriftrollen lagen am Boden umher, dort drüber war sogar ein Regal umgestürzt, und jede Oberfläche war überzogen mit einer dicken Schicht Staub, die Kjalls Niesreiz ganz gewaltig reizte. Offenbar war schon seit Monden niemand mehr hier gewesen und hatte sich um die Ordnung gekümmert.

Kjall versteckte seine Holzkiste mit den Berichten zu verschiedenen Zwergenartefakten unter einem Lesepult und überlegte. Diese letzte Begegnung mit Folla und dem Feerich Fanatos hatte ihn etwas aus der Bahn geworfen. Warum hatte Folla ihn erkannt, ehe er in die Vergangenheit gereist war? War eine solche Interaktion über verschiedene Zeitlinien hinweg möglich? Oder war die Zeitlinie doch vorgeschrieben und Andors Entstehung unvermeidlich? Konnte man dies überhaupt überprüfen?

Kjall lauschte an der Tür, die die Schwarzen Archive mit dem Rest des Baums der Lieder verband. Keine Stimmen waren zu hören. Schliefen die Bewahrer allesamt?

Die Tür erwies sich als hartnäckigeres Hindernis als erwartet, zumal Kjall seinen Casamatuc in der Vergangenheit liegen gelassen hatte. Beim Netz der Schicksalssporne, er würde nicht in diesem Raum verrotten! Endlich, mit etwas Gewalt und einem eisernen Kerzenhalter, schaffte Kjall es, das Schloss aufzubrechen. Die Tür schwang mit einem lauten Knarren auf und Kjall duckte sich zurück ins Dunkel, für den Fall, dass das Geräusch einige Bewahrer geweckt hatte. Nichts. Nach einigen Minuten brachte Kjall genug Mut auf, die Schwarzen Archive zu verlassen. Er huschte aus dem Türrahmen und blickte zurück auf die aufgebrochene Tür. Jemand hatte in großen schwarzen Lettern „Betreten verboten per Dekret des \st{Statthalters Ken Dorr} Obersten Bewahrers“ darüber gepinselt.

Zumindest sah der restliche Teil des Baums der Lieder nicht so verlassen aus wie die Schwarzen Archive. Der Boden war sauber, die Türen waren beschildert und drei Stockwerke weiter oben sah Kjall sogar noch ein Lichtlein brennen – wohl irgendein Mitglied des Bewahrerordens, dessen Enthusiasmus für alte Schriften keine Tageszeit kannte. Ein Gefühl, von dem Kjall selbst nur allzu gut wusste, wie es sich anfühlte.

Kjall wusste ganz genau, wo er hinwollte: Im zweiten Stock, in der großen Bibliothek, gab es eine Sammlung von Aufzeichnungen zu den vier mächtigen Schilden aus uralter Zeit, in welcher Kjall bereits geschmökert hatte. Zu manchen war erheblich mehr bekannt als zu anderen, Zum Feuerschild hatten die Bewahrer aus Kjalls Zeit etwa bloß eine kleine, mickrige Erwähnung aus der Zeit der Trollkriege gefunden. Oh, der mächtige Feuerschild! Kjall hatte seine Kraft erleben dürfen, als sich ihm in der Vergangenheit die Helden von Andor und dieser ach so ehrenvolle Schwertmeister Harthalt entgegengestellt hatten. Es gab keinen Zweifel daran, dass der magische Schild, den Harthalt geführt hatte, der Feuerschild gewesen war! Wie dieser Lackaffe ihn wohl errungen hatte? Und ob er überhaupt wusste, was für einen Schatz er da trug? Wahrscheinlich nicht, die Schildzwerge hatten die Geschichte des Feuerschilds gut unter Schutz gehalten. Nun gut, es brauchte Kjall nicht zu kümmern. Ein hämisches Grinsen zog sich über sein Gesicht: Vielleicht könnte er Harthalt sogar mithilfe der Sphäre einen Besuch abstatten gehen und ihm den Feuerschild abknöpfen. Das wäre doch was: Nicht nur einen, sondern gleich zwei mächtige Schilde zu führen. Kjall wurde ganz hibbelig bei der Vorstellung. Fokus!

Fokus! Jetzt ging es zunächst einmal darum, herauszufinden, ob Kjall durch seinen Diebstahl des Silberschilds wirklich die Vergangenheit verändert hatte. Die große Bibliothek war einfach genug zu finden, und bei dieser Beschilderung fand Kjall die große Sammlung zu den vier mächtigen Schilden noch viel einfacher. Wo war denn die Chronik? Ah, da! Silberschild, Silberschild, Silberschild... nichts!? Wie sah es unter ‚Sturmschild‘ aus? Ah, hier! „lagerte für hunderte von Jahren in der Waffenkammer der Schildzwerge, ehe er von Fürst Hallwort dem Großen nach Silberhall befördert wurde.“

„Feuer und Spucke“, murmelte Kjall, „Feuer und Spucke...“

Es war, als hätte er den Schild nie gestohlen. Hatte er die Vergangenheit gar nicht verändert? Wobei, vielleicht war das hier auch nur eine andere Zeitlinie. Zum Glück konnte Kjall zumindest das überprüfen. Er raste in eine andere Sektion der großen Bibliothek und suchte die Chronik zu den Unruhen am Baum der Lieder. Zur Zeit der Trollkriege... in einem Herbst... das Jahr, das Jahr... hier war es: „Jork (?), rothaariger Zwerg, Einbruch in die Schwarzen Archive, Diebstahl zahlreicher Schriften, Mord. Verfügt über unbekannte Portaltechnologie. Zwei Unruhen in einem Abstand von rund einem Mond. Beziehung zu Thogger? Entflohen aus den Kerkern der Schildzwerge gemäß Aussage von Wächter Bort.“

Sein Besuche am Baum der Lieder waren niedergeschrieben worden. Das machte es eindeutig. Dies war dieselbe Zeitlinie, in welcher Kjall den Silberschild gestohlen hatte. Und dennoch hatte er keine Veränderung am Verlauf der Geschichte ausgelöst. Kein Paradoxon hatte ihn verschlungen. Hatte Fürst Hallwort mit seinen Behauptungen recht gehabt? Standen das Vergangene und die Zukunft in Stein gemeißelt? Das würde doch bedeuten, dass irgendjemand irgendwann den Silberschild, den er immer noch auf seinem Rücken trug, zurück in seine Zeit bringen würde. Hieß dass, das Kjall den Schild wieder verlieren würde, verlieren musste?

Kjall sank neben dem Pult zu Boden. Wenn diese Überlegungen stimmten, so konnte er tun und lassen, was er wollte, und dennoch nicht verhindern, dass Hallwort in den Norden zog. Nicht verhindern, dass Hallgard Xoll in seinen Tod schickte. Nicht verhindern, dass Kram zum Fürsten ernannt wurde. Nicht verhindern, dass Radan verbannt wurde. Eine Schande war das! Wenn die Geschichte feststand, konnte Kjall geradesogut gar nicht erst versuchen, das Schicksal zum Guten zu wenden. Aber... Moment mal, er konnte das Ende der Geschichte herausfinden. Ja, er konnte die Zukunft Caverns erfahren, hier, in diesem Augenblick! Hier in dieser Bibliothek lagerten die Aufzeichnungen aller Jahrhunderte, die zwischen Kjalls Zeit und dieser Zukunft liegen mögen, und Kjall konnte sie studieren. Wobei, wenn er darüber nachdachte, dann gab es sogar noch einen schnelleren Weg. Einen visuelleren Weg.

Von neuem Tatendrang erfüllt, sprang Kjall auf. Er verließ die große Bibliothek und hetzte die Wendeltreppe mit den viel zu hohen Stufen hoch. Er wollte ganz oben an die Spitze. Dort, am höchsten Punkt des Baums der Lieder, befand sich eine kleine Aussichtsplattform, von der aus man den gesamten Wachsamen Wald und das Rietland überblicken konnte – und natürlich die Ausläufer des Grauen Gebirges, unter denen ein Teil des Zwergenreichs Cavern lag.\bigskip



Kühle Winterluft schlug Kjall entgegen, als er das Ende des Wendeltreppe erreichte und die Tür zur Aussichtsplattform des Baums der Lieder erreichte. Er rutschte beinahe auf der dünnen Schneeschicht aus, die sich hier angelagert hatte, konnte sich gerade noch fassen und rannte dann nach vorne zur Balustrade, wo er sich auf die Zehenspitzen stellte, um knapp darüberzusehen. Der Silberschild fiel klappernd zu Boden. Kjall riss ihn an sich und bemerkte kaum, dass der kalte Wind um ihn herum sich legte und Wärme in seine Glieder strömte. Er war zu beeindruckt von der Aussicht, die sich ihm zeigte. Beeindruckt – und erschrocken.\bigskip



Die Nacht lag über Andor, nur der rote Mond und einige Sterne standen am Himmel. Dennoch erglühte das Rietland in einer Lichterpracht. In der Ferne leuchteten die Fenster der Rietburg, dieses verhassten Gebäudes, als säße im Innern der Häuser und Türme immer noch der Tag gefangen. Seltsame, golden glitzernde Seile führten von den Türmen und Zinnen der Burg zu Masten und von da ins weitere Rietland... aber wo war denn das Rietland? Wo früher nichts als trockenes Gras im Wind geweht hatte, standen nun Häuser dicht an Häusern, von der Rietburg bis hin zum Südlichen Wald. Der südliche Wald hatte sogar etwas an Größe eingebüßt, um Platz für weitere Siedlungen zu lassen. Das waren keine einfachen Bauernkaten, jedes einzelne dieser Häuser ragte mehrstöckig in den Himmeln und hatte Anbauten, die die Taverne zum Trunkenen Troll vor Neid erblassen lassen würden. Noch dazu schien bei all diesen Häusern Licht aus den Fenstern hinaus, als hielte man Gevatter Tag höchstpersönlich gefangen! Es war absurd. Wo einst der freie Markt gelegen hatte, stand nun ein breiter Turm – zwergische Machart, das erkannte Kjall sofort. Immerhin schien die Taverne zum Trunkenen Troll, dieser winzige Punkt am Rande des Lichtermeeres, noch so auszusehen, wie Kjall sie kannte – aber auch zu diesem Gebäude führten eigenartige hängende Seile, an denen hin und wieder ein bläulicher Blitz entlangraste, und auch aus den Fenstern der Taverne glühte goldenes Tageslicht, wie es selbst die Fackeln von Cavern nicht zustande brachten.

Kjall wurde Angst und Bange, als ihm bewusst wurde, wie viele Menschen in einem einzelnen Haus leben konnten. Diese Unmengen an Behausungen im Rietland waren bestimmt noch nichts im Vergleich mit der schieren Anzahl an Menschen, die sich dort versammelt haben mussten und taten, was Menschen eben so tun.

Mit Schaudern erkannte Kjall, dass selbst das Fahle Gebirge, dessen Gipfel immer in den Wolken zu stecken schienen, von diesen seltsamen, golden leuchtenden Seilen überspannt war... und waren das etwa Häuser? Häuser auf dem Fahlen Gebirge?! Kjall rieb sich die Augen. Jawohl, da schienen tatsächlich noch weitere Hütten zu stehen. Der warme Schein aus ihren Fenstern beleuchtete den dunklen Smog, der aus ihren Schornsteinen dampfte und den Hang des Fahlen Gebirges entlang zog um sich in einer Schwarzen Wolke am Gipfel zu sammeln, die sich mit dem allgegenwärtigen Weiß des Schnees biss.

Das Gebiet südlich der Rietstadt und des Südlichen Waldes lag im Dunkeln der Nacht. Immerhin hatten die Menschen nicht auch dort noch Häuser hingepflanzt. Vermutlich mussten sie Felder anlegen. Kjall wurde schwindelig, als er sich überlegte, was dieser Moloch einer Stadt für Nahrungsbedürfnisse haben musste. Sein Blick schweifte zur Narne, auf der er... waren das Boote? Das konnten doch unmöglich Boote sein, die waren ja viel zu groß! Und was war dieser golden leuchtende Dampf, den sie ausstießen? Konnte es sein, dass sie von derselben Technologie angetrieben wurde, die diese goldenen Leitungen im Rietland verteilten?

Kjalls Blick folgte dem Lauf der Narne in Richtung des Hadrischen Meeres, und erstarrte erneut. Die Nebelinseln, so wenig er von ihnen auf diese Distanz auch erkennen konnte, waren ebenfalls von farbigem Licht erfüllt, das den Nebel und kleine Teile des umliegenden Meeres erglühen ließen. Ein weiteres Glühen erhellte eine unförmige dunkle Erhebung im Hadrischen Meer in der Mitte zwischen Sidra und Silberland. Kjall konnte sie nicht näher einordnen. Es sah fast aus, als hätte man einen verdammten Riesenkraken versteinert.

Kjall senkte seinen Blick auf den Wachsamen Wald, und auch hier, zwischen den immergrünen Bäumen, waren hier und da golden glänzende Leitungen gespannt worden, die immer wieder kurz aufglühten und dann wieder dunkler wurden. Immerhin lag der Wald abgesehen davon im Dunkeln. Die Bewahrer und Dorfbewohner schienen zumindest die Ruhe der Nacht zu achten.

Einen letzten Ort in diesem Panorama gab es, den Kjall noch nicht betrachtet hatte, weil er sich davor fürchtete, was er sehen würde. Cavern, die Heimat der Schildzwerge. Sein Heimat. Wie hatte das Zwergenreich den Zahn der Zeit überstanden? Kjall schlich vorsichtig zum anderen Ende der Balustrade und hievte sich hoch.

Das Graue Gebirge stand schwarz vor dem Sternenhimmel, und keinerlei goldene Leitungen überzogen die Berge. Licht brach dennoch aus ihm hervor, aus dutzenden Türmen und Turmgruppen, die scheinbar zufällig über die Berge verteilt waren. Meisterwerke der zwergischen Handwerkskunst, die zu Kjalls Zeiten noch nicht existiert hatten. Selbst der Verlassene Turm strahlte aus seinen Fenstern, und die bröckelige Treppe, die vom Fuße des Gebirges zum Turm führte, wirkte blitzblank, so gut wie neu geschaffen! Der südliche Eingang nach Cavern war besonders hell erleuchtet, und nicht nur im erwarteten Gelb, nein, da schwebten grüne, blaue und rote Sphären voller Licht, die den Eingang nach Cavern beleuchteten und die Dutzenden von Zwergen beschienen, die sich dort aufhielten. Sie bewegten sich rhythmisch und schienen eine Art Zeremonie abzuhalten. Der Wind trug von Zeit zu Zeit kurze Fetzen ihres Gegröles an Kjalls Ohr. Dessen Aufmerksamkeit schweifte aber schnell wieder zurück zu den eleganten Bauten, die das Graue Gebirge überzogen. Wenn es schon von außen so aussah, wie musste es dann erst innen sein? Oder auf der anderen Seite des Grauen Gebirges, das hier waren ja sogar nur die nördlichen Ausläufer!

Es gab keinen Zweifel: Das Cavern war aufs Neue erblüht und erstarkt, trotz aller Widrigkeiten, die sich ihm in den Weg gestellt hatten. Kjall konnte kaum genug kriegen vom Anblick der leuchtenden zwergischen Handwerkskunst unter dem Sternenhimmel. Er wischte sich eine Träne aus dem Auge. Die Zukunft Caverns... sie war gut.

\newpage
\section{Gestohlen}

Kjall wusste nicht, wie lange er dagestanden und die Pracht des erstarkten Caverns bestaunt hatte, als ihn ein nicht allzu ferner Ruf eines Skrals zurück in die Realität holte. Genauer gesagt wurde ihm plötzlich bewusst, dass er nicht als einziger auf der Aussichtsplattform oben am Baum der Lieder saß und die Umgebung begutachtete. Eine junge Bewahrerin in grauer Kleidung hatte sich zu ihm gesellt und musterte ihn neugierig.

„Zum Gruße“, setzte Kjall vorsichtig an, „Ich... bin Jork.“

„Ich bin Josella“, antwortete die Bewahrerin mit einem freundlichen Lächeln, „und es freut mich, dass Ihr zu uns gefunden habt. Dieser Ort ist wahrscheinlich mein Lieblingsplatz im gesamten Wachsamen Wald. Hier kann man viel Zeit verbringen und die Umgebung begutachten. In der Nacht ist es besonders schön. All diese Lichter...“

Kjall nickte. Die Aussicht war wirklich atemberaubend bezaubernd, wenn auch für ihn wohl nicht aus denselben Gründen wie für sie.

„Wie seid Ihr an den Wachen unten am Baum der Lieder vorbeigekommen?“, fragte Josella nun. Reine Neugierde sprach aus ihrem Gesicht, keine Feindseligkeit.

Unverfrorene Lügen hatten ihm noch selten geschadet, also meinte Kjall fröhlich: „Ich habe mit den Wachen einfach geredet, natürlich. Ihnen ist mein Abliegen sehr verständlich gewesen.“

„Na gut, ich beiße“, lächelte Josella, „Wollen Sie mir Ihr Anliegen verraten? Was suchen Sie des nachts oben auf einer Aussichtsplattform, welche eigentlich einzig und allein uns Bewahrern zugelassen wäre?“

„Nun, das ist leicht. Ich wollte das glorreiche Cavern bei Nacht erblicken. Ich hatte einst ein Gemälde gesehen, welches dieses Blick zeigte, ein Gemälde, welches mich sehr faszinierte. Bis eben konnte ich nicht glauben, dass es wirklich so schön aussieht.“

Josella nickte. Stille breite sich zwischen den beiden aus, wie sie gemeinsam an der Balustrade standen und das Graue Gebirge betrachteten. Der starker Duft von Wolfskraut drang an Kjalls verfrorene Nase. Hatte Josella welches zerrieben, ehe sie hier nach oben gekommen war? Kjall musste aufpassen, dass es ihm nicht die Sinne vernebelte. Die Bewahrerin schien weiterhin freundlich gesinnt zu sein, auch wenn Kjall sich sicher war, dass sie ihm seine Geschichte nicht abkaufte.

„Darf... darf ich Euch etwas fragen?“, fragte Kjall.

Josella nickte erwartungsvoll.

„Kennt Ihr Euch gut aus mit der Geschichte des Zwergenreichs? Genauer gesagt der Zeit unter der Führung von Fürst Kram?“

Josella druckste etwas herum: „Huch, das ist lange her... lange her, dass es geschehen ist, und auch schon lange her, dass ich die Geschichte Caverns gebüffelt hatte.“

Kjall blieb still und wartete darauf, dass Josella ausführte.

„Nun, gut, wenn Ihr unbedingt meine Kenntnisse überprüfen wollt... Fürst Kram herrschte über Cavern zwischen den Jahren 72 und 76 nach andorischer Zeitrechnung, die kürzeste Herrschaftsperiode aller Fürsten, die seither hier geherrscht hatten.“

Kjall horchte auf.

„Ich glaube nicht, dass ich näher auf Fürst Krams überraschenden Tod eingehen muss, diese Anekdote sollte inzwischen jedem bekannt sein...“

„Tut so, als hätte ich noch nie davon gehört.“

Josella warf ihm einen verwirrten Blick zu und meinte dann: „Nun... nun... Fürst Radan hatte ja bereits einmal versucht, einen Aufstand gegen Fürst Kram anzuzetteln, kurz nach der Krönung Fürst Krams, und war dafür verbannt worden. Nachdem sich Radan allerdings in der Schlacht gegen Krahd auf die Seite der Zwerge entschieden hatte, wurde dessen Verbannung aufgehoben und Radan mit aller Ehre zurück nach Cavern eingeladen.“

Kjall spitzte die Ohren und grinste innerlich schon über beide selbige hinaus.

Mit gequältem Gesichtsausdruck fuhr Josella fort: „Kaum drei Monde nach der Rückkehr aus dem Lande Krahd startete Radan einen weiteren Aufstand, diesmal erfolgreich. Fürst Kram wurde im Bad überrascht und hinterrücks... nun ja... Fürst Radan wurde Herrscher über Cavern und regierte mit strenger Faust, stets die Zukunft im Blick behaltend. So expandierte er Cavern, brachte alte Zwergenfesten wieder zum Aufbau und befahl die Errichtung neuer Meisterwerke. Den Kontakt zu Andor und dem Barbarenland – und den Zwergen aus den Silberminen, die ihm die Krone absprachen – brach er größtenteils ab, weswegen wir leider nur über wenige Berichte aus dieser Zeit verfügen. Aber es lässt sich nicht leugnen, dass Cavern unter seiner Führung zu alter Stärke zurückfand. Soll ich fortfahren?“

„Ich bitte darum“, nickte Kjall eifrig. Radan als Fürst, das hätte er sich nie träumen lassen!

Josella schien sich etwas zu entspannen, als sie fortfuhr: „Radans einziger Sohn und Nachfolger, Fürst Breggo, stärkte die Bunde Caverns mit seinen Nachbarn wieder, während er stets auf dessen Unabhängigkeit achtete, und dass bei den unzähligen Händel mit der Außenwelt Cavern stets den größten Profit zulegte. Hatte man sich zur Zeit von Radans Führung noch manchmal vor kriegierischen Auseinandersetzungen zwischen Cavern und den Reichen der Menschen gefürchtet, so konnte das gesamte Land unter Breggos Führung aufatmen.“

Kjall knirschte mit den Zähnen. Er hätte nur zu gerne gesehen, wie die Heere der Menschen unter dem Ansturm Caverns zurückgedrängt wurde. Aber es konnten ja nicht alle Wünsche in Erfüllung gehen.

„Und wer herrscht jetzt über das Zwergenreich?“

Josella runzelte ihre Stirn leicht, verlor aber nicht ihr Lächeln: „Woher kommt Ihr denn, dass Ihr nicht einmal wisst, wer... verzeiht, das war unhöflich. Die regierende Herrscherin ist Meria, Fürstin unter dem Berge und Trägerin der Schildkrone, Tochter des Breggo, Sohn des Radan, Sohn des Eidin. Sie ist jung, erst um die 90 Jahre alt. Ich habe den Zeitpunkt ihrer Krönung miterlebt, da war ich noch ein Kind. Damit meine ich, dass ich von ihrer Krönung gehört habe, als sie abgehandelt war. Menschen waren zu diesem festlichen Anlass natürlich keine eingeladen.“

Kjall nickte hocherfreut. Er mochte den Verlauf der Zeit nicht so steuern können, wie er es gedacht hatte, aber er konnte sich glücklich schätzen, dass die Zukunft Caverns so rosig aussah. Nun hatte er genug gesehen von dieser seltsamen Zukunft voller überfüllter Menschenhäuser und goldener Leitungen. Kjall konnte sich nun entspannen, zurück in seine eigene Zeit reisen und... und ja, was tun?

Was wollte Kjall denn überhaupt in seiner eigenen Zeit? Klar, es wäre schön, Radans Aufstieg hautnah mitzuerleben, aber noch schöner wäre es zum Beispiel, den Urzeiten Caverns einen Besuch abzustatten, vielleicht sogar an ihnen beteiligt zu sein. Mithilfe der Sphäre stand Kjall die gesamte Zeitlinie Andors offen, und nun wusste er auch mit ziemlicher Sicherheit, dass er nichts tun würde, was der Entwicklung Caverns schaden könnte – vielleicht war er sogar an deren Entstehung beteiligt. Die Blütezeit der Drachen, als es noch keine elende Burg im Rietland und keine Krahder in Krahd, ja sogar keine Bewahrer im Wald gab! Als Drachen durch die Lüfte flogen, im Bündnis mit den Zwergen, die die Tradition noch zu achten wussten. Dorthin wollte Kjall. Und dorthin würde er auch kommen.

Aber noch nicht sofort. Kjall betrachtete den Silberschild in seiner Hand und ein alter Instinkt erwachte wieder in ihm. Ein Instinkt, den er lange Zeit unterdrückt hatte. Der Sammlertrieb! Kjalls Lebenserwartung war noch lange. Solange die Sphäre genug Saft enthielt, konnte Kjall nach Belieben durch die Zeit hopsen. Und neuen Treibstoff herzustellen war auch nicht die unüberwindbare Aufgabe, für die er sie einst gehalten hatte. Flüssiges Feuer, Tiefsand vom Geheimen See... das konnte er alles hier in der Nähe finden. Die Zeit, die Herstellung zu perfektionieren, hatte Kjall ja. Die Urzeiten Caverns würden auf ihn warten, eines Tages. Zunächst einmal galt, es seine Sammlung vervollständigen. Die vier mächtigen Schilde. Ein Hadrischer Kompass. Die magischen Waffen. Kjalls Herz schlug immer schneller, als er sich erst so richtig dessen bewusst wurde, was für Möglichkeiten er nun hatte. Die ganze Welt stand ihm offen!\bigskip



„Seid Ihr noch bei uns, Jork?“, fragte Josella vorsichtig und stupste Kjalls Schulter. Kjall schüttelte sich, sein Blick klärte sich, und ihm wurde wieder bewusst, wo er sich befand.

„Natürlich, Josella. Aber nicht mehr lange, denn ich muss mich gleich wieder auf die Socken machen.“

Josella senkte den Kopf zur Verabschiedung: „Wenn Ihr meint... seid Ihr sicher, dass Ihr euch im dunklen Wald zurechtfinden werdet? Es streifen noch immer Mitglieder dieser Trollhorde aus dem Barbarenland durch die Gegend, die unseren Jägern wie durch Zauberhand durch die Maschen zu schlüpfen schienen. Wir habe einen Experten aus Werftheim gerufen, aber auch der meinte, wir müssten einen Zauberer... und weg ist er.“

Weg war er tatsächlich, denn während Josellas Monolog hatte Kjall sich von der Brüstung entfernt und war wieder im Innern des Baumes verschwunden, wo man ihn rasch die Treppe runterklappern hörte.

Ich lauschte Kjalls leiser werdenden Schritten und trat dann aus dem Schatten hervor.

„Ich muss dir von ganzem Herzen danken, Josella. Du weißt nicht, wie sehr du gerade unseren beiden Reichen geholfen hast.“

Josella blickte mit zusammengekniffenen Augen auf mich herunter. Jetzt lächelte sie nicht mehr: „Das hat sich falsch angefühlt. Und das war es auch. Wir Bewahrer vom Baum der Lieder sind der Tradierung der Wahrheit verpflichtet, nichts als der Wahrheit. Was bringt es, ihn jetzt zu belügen? Er wird es doch ohnehin herausfinden.“

Ich schüttelte meinen Kopf und sprach durch den Raureif in meinem Bart: „Du hast ihm das gegeben, was er wollte. Du hast ihm einen letzten Grund gegeben, die Zukunft Zukunft sein zu lassen und sich anderen, vergangenen Dingen zuzuwenden. Dingen, die er erwiesenermaßen nicht zerstören kann. Du hast ihm und unseren Ländern zu Frieden verholfen. Er wird sich wohl kaum die Mühe machen, deine Geschichte nachzuprüfen. Und so wird er nie auf den Gedanken kommen, sie doch noch zu den Ungunsten von Krams Blutlinie modifizieren zu wollen.“

Josella schien noch nicht überzeugt, bohrte aber auch nicht weiter nach. Stattdessen fragte sie: „Er trug einen Eurer mächtigen Schilde. Warum habt Ihr ihn ihm nicht einfach abgenommen?“

Ich grinste: „Um dasselbe gleich noch drei weitere Male tun zu müssen? Ich warte lieber, bis er alle vier gesammelt hat und greife sie mir dann.“

„Aber... aber wenn Ihr ihn heute geschnappt hättet, dann hätte er ja gar nicht alle vier...“

„Manchmal ist es besser, nicht zu sehr darüber nachzudenken. Ich weiß, dass er eines Tages alle vier finden und vereinen wird. Das ist das, was zählt.“

Josella schüttelte wieder ihren Kopf: „Wisst Ihr überhaupt, worauf Ihr Euch da einlässt?! Das sind die vier mächtigen Schilde aus der Urzeit! Wie könnt selbst Ihr auch nur davon träumen, gegen einen Träger der Vier anzukommen?“

Ich grinste: „Dieser Kjall ist ein kleiner Bengel, der keine Ahnung hat, mit welchen Mächten er da spielt.“

Ich wuchtete den schweren Sack von meinem Rücken, wo er mit einem Scheppern auf den Holzboden der Aussichtsplattform fiel. Ich öffnete den Sack und ließ Josella einen Blick hineinwerfen.

Meine Grinsen wurde breiter: „Davon abgesehen besitze ich die Vier doch auch.“

Josella machte große Augen und war dann still.\bigskip




\az{Jahr 72}

Ich war zu früh. Typisch.

Die Reise in und durch die Tiefminen war überraschend ereignislos verlaufen, auch wenn mich nur der geschickte Einsatz des Sternenschilds davon abgehalten hatte, gleich bei meinem Eintritt in die Tiefminen durch einen abrupten Feuerstoß gebrutzelt zu werden. Ich hätte vielleicht doch lieber den Zugang durch die alte Zwergentür am Baum der Lieder nehmen sollen, aber da ich es so sehr mochte, durch diese altbekannten Stollen zu spazieren und den Stein unter meinen Fingerspitzen zu spüren, wagte ich den Umweg über die östlichen Tiefminen.

Still und ruhig lag es da, das zerfallene Museum Kjalls. Eine große Zwergenaxt lag vor dem Eingang, aber sie war das einzige Artefakt, welches hier noch vollständig war. Das Innere des Museums war von einer großen Staubschicht überzogen, Steinbrocken lagen kreuz und quer über den Raum verteilt, durchmischt mit verkohlten Holzsplittern und verbogenen Metallplatten. Ein trauriger Anblick. Was auch immer hier gewütet hatte, es war eine äußerst destruktive Macht gewesen.

Die Kalibration für einen Zeitsprung von über einem Jahrhundert war äußerst komplex, und ich wollte lieber zu früh als zu spät auftauchen. So kam es, dass ich noch einen halben Tag im Jahr 72 zu warten hatte, ehe Kjall zur Tür hineinspazierte. Stunden, die ich unter anderem damit verbrachte, Kjalls Museum in seiner Blütezeit zu bestaunen.

Kjall hatte bereits vor dem Erringen der Sphäre eine beeindruckende Sammlung uralter Zwergenartefakte erreicht, aber nun war die schiere Anzahl von Ausstellungsstücken einfach absurd. Eine andorische Flöte, die Brandur selbst gespielt hatte. Ein hadrischer Kompass, der den Seekrieger Ruuf damals zur mythischen Insel Danwar geleitet hatte. Boords Hammer, mit dem er die große Statue in der Halle der vier Schilde aus dem Stein gehauen hatte. Selbst einer von Nehals Fangzähnen! Alle Artefakte waren penibel beschriftet, von Staub befreit, auf verschiedensten Sockeln zur Schau gestellt. Ich kam kaum aus dem Staunen heraus.

Das Hauptaugenmerk des geheimen Museums waren natürlich die Sockel der vier mächtigen Schilde aus uralter Zeit. Drei hatte Kjall bereits erbeutet, und sie standen auf Hochglanz poliert nebeneinander: Der Sternenschild, der Bruderschild, und der Sturmschild. Mein Herz pochte schneller, als ich sanft über ihre glatten Oberflächen strich und sie unter meiner Berührung leise aufsummten, als würden sie mich erkennen. Meine Meisterwerke. Nur der Feuerschild fehlte noch, stattdessen war auf seinem Sockel eine Skizze des Schwertmeisters Harthalt abgebildet, wie er den markanten dreieckigen Schild führte.

Der Sockel des Feuerschilds war nicht der einzige, der mit einer Skizze statt eines Artefakts belegt war. Wahrscheinlich erfuhr Kjall schneller von neuen Artefakten, als er die bekannten einsacken konnte, sodass es ihm gar nie möglich sein würde, eine vollständige Sammlung zu besitzen. Orweyns Hammer der Stärke und Varatans Helm der Macht hatte Kjall noch nicht errungen, ebenso fehlte ihm ein Amulett des roten Mondes und ein Hadrischer Spiegel. An ein vollständiges Hadrisches Stundenglas war Kjall offenbar auch noch nicht gekommen, der entsprechende Sockel war nur mit einigen Scherben, etwas Sand und einem halben Holzgestell belegt. Dafür bemerkte ich mit Staunen, dass Kjall eine der drei magischen Waffen aus Hadria errungen hatte! Varlion, das Flammenschwert, steckte in seiner Scheide in der hinteren linken Ecke des Museums. Bei den Hörnern des Urtrolls, wie war Kjall daran gekommen? Hatte er es der Zauberin Eara abgeknöpft, nachdem diese damit durch das Schwarze Portal nach Andor gereist war? Ich überlegte mir kurz, das Schwert einzustecken, entschied mich dann aber dagegen. Ich war wegen meiner Schilde hier, da musste ich mir nicht noch mehr Aufgaben aufbürden.\bigskip



Als Kjall endlich ins Museum trat, hatte ich es mir in einer Ecke des Museums gemütlich gemacht und mir die Kapuze des Hral übergeworfen, auf dass ich beinahe eins mit dem Schatten der Ecke wurde. Kjalls Kleidung war angekokelt und er murmelte fluchend etwas von einem elenden Sektenanführer, der diesen Dunklen Tempel ohnehin nicht verdient hatte. Dabei verrieten Kjalls blitzende Augen seine eigentlich prächtige Laune. Er übersah mich vollständig, ganz wie ich es erwartet hatte, und trat mit einem breiten Grinsen vor die Sockel der vier Schilde. In seiner linken Hand sah ich seine Sphäre aufblitzen, bevor er sie in seine Tasche wandern ließ. In seiner rechten Hand hielt er den Feuerschild, welchen er triumphierend auf seinem Sockel platzierte. Dann trat er zurück und stand einfach still da.

Ich hielt mich ebenfalls still und wartete auf eine Reaktion.

Kjall setzte sich andächtig auf den Boden und betrachtete seine vier Schätze. Ich hätte schwören könne, aus seinem Augenwinkel eine Träne laufen zu sehen. Dann war der Bann gebrochen und Kjall lachte auf, klatschte in seine Hände und sprang, nein, tanzte förmlich um die Sockel der vier Schilde herum. Ich ließ ihn einige Minuten des Triumphs genießen. Dann aber hielt ich es für angebracht, die Kapuze des Hral abzustreifen und in Kjalls Applaus einzustimmen.

Wie erwartet trat Kjall erschrocken zurück und stolperte beinahe über den Sockel mit dem frisch errungenen Feuerschild. Er fasste sich, richtete sich stolz auf und rief: „Wer... wie...“

Kjall blickte die vier Schilde an und schien zu verstehen: „Kreatok?!“

Ich nickte.

Kjall griff sich an seine gefurchte Stirn und meinte: „Drachenfeuer und Zwergenspucke, wie hast du mich gefunden?“

Ich hielt einige Fetzen uralten Pergaments in die Höhe, die mit einer sehr kleinen Handschrift überzogen waren: „Eine mysteriöse Nachricht aus der Vergangenheit. So relevant ist das nicht. Wichtiger ist, was du hier für eine wundervolle Sammlung angelegt hast. Ich habe mich etwas umgesehen, und ich bin echt beeindruckt!“

Kjalls Schultern sackten herunter und ein erleichterter Ausdruck flog über seine Züge: „Meister Kreatok, es ist mir eine unglaubliche Ehre... es gibt so viel, was ich dich fragen könnte. Wie... wie funktionieren diese Schilde? Könnte man weitere produzieren? Was hat dich dazu geritten, den Unterirdischen Krieg anzuzetteln? Was für Abenteuer hast du mit der Sphäre erlebt? Könnten wir die Vergangenheit irgendwie ändern?“

Ich gebot Kjall, seinen Redeschwall einzustellen. Mein Vater hatte mich gelehrt, schlechte Neuigkeiten nicht unnötig lange hinauszuzögen, also fuhr ich fort: „So sehr mir dieses Museum auch gefällt, so muss dir in der Zwischenzeit bestimmt bewusst geworden sein, dass diese Sammlung hier nicht auf ewig Bestand haben können wird. Es ist unklug, derart mächtige Gegenstände so nahe beieinander zu halten, und erst recht unklug, wenn eine solche Sphäre dabei ist. Falls jemand darauf stoßen sollte, der sich nicht so sehr aufs verborgene Vorgehen versteht wie du, Kjall, so wären die Folgen undenkbar. Jemand sollte diese Schilde und die Sphäre wieder an ihre rechtmäßigen Orte zurückbringen. Und dieser jemand bin ich.“

Mit diesen Worten trat ich vollständig aus dem Schatten und Kjall sog scharf Luft ein, als er meine Ausrüstung sah. Ich trug meinen Bruderschild in meiner linken und meinen Sturmschild in meiner rechten Hand. Meinen Sternenschild hatte ich mir auf den Rücken geschnürt, wo er strahlend hell zu leuchten begann und das geheime Museum Kjalls in einen türkisen Schein tauchte. Meinen Feuerschild hatte ich mir auf die Brust gebunden. Er war der einzige Schild, von dem ich mir geschworen hatte, ihn heute nicht einzusetzen.

Eine dunkle Macht lag über diesem Schild, eine dunkle Macht, die den Unterirdischen Krieg und den Zerfall des Zwergenreichs ausgelöst hatte. In meiner eigenen Zeit hatte ich diesen Schild noch gar nicht geschmiedet, aber ich wusste, dass ich ihn eines Tages schmieden würde, und ich wusste, dass ich danach den silbernen Feuerschlund entfesseln würde. Und Nehal würde brennen. Mein Nehal. Dafür würde es eine Erklärung geben müssen, da war ich mir sicher. Aber ich kannte sie noch nicht. Würden mir die ganzen Zeitreisen den Verstand verdreht haben? Würde mich ein Fluch befallen haben? Ich war ratlos, und ich fürchtete mich davor, die Wahrheit zu erfahren. Aber im Moment musste ich mich nicht darum kümmern. Im Moment zählte nur, dass ich nicht aus Versehen die finstere Macht des Schildes entfesselte.

„Ich hoffe, dass wir diese Angelegenheit friedlich klären können“, sprach ich weiter, „Wenn du mir die Sphäre und die Schilde übergibst, bringe ich sie weg von hier und kehre dann zurück, um mich in aller Ausführlichkeit mit dir über die Wirkungsweise der Schilde zu unterhalten.“

Ein wirrer Blick trat in Kjalls Augen: „Kreatok, deine Zeit ist vorüber. Du hast nie das volle Potential deiner Sphäre erkannt, oder sie nicht ausgiebig genug genutzt. Diese unermessliche Macht... ich könnte zum Todeszeitpunkt meines Vaters reisen und ihn in die Zukunft holen. Ich könnte ein Bündel Sternkraut aus dem Grauen Gebirge holen und es Thogger vor die Füße werfen, ehe er überhaupt in den Süden aufbräche. Ich könnte tun und lassen, was ich will, und ich tue es dennoch nicht!“

„...weil du weißt, dass du damit nie eine Veränderung erzielen würdest...“

„Nein, weil ich mich beherrschen kann! Mit dieser Sphäre bin ich ein Gott, und doch halte ich mich immer noch im Verborgenen, weil mir dieser große Rummel nichts bedeutet! Dieses Museum schadet niemandem, und niemand außer mir wird es je zu sehen bekommen. So lasse mir doch diese kleine Freude!“

„Ich frage nur nach der Sphäre und den Schilden, den restlichen Kram mögest du behalten.“

„Unterschätze mich nicht, Kreatok!“, fauchte Kjall nun, „Ich besitze dieselbe Macht, dieselben Artefakte wie du, aber im Gegensatz zu dir habe ich sie sinnvoll eingesetzt. Ich habe hart gearbeitet für diese Sammlung...“

Kjall schüttelte enttäuscht seinen Kopf.

„Ich hätte gedacht, dass wenigstens du erkennen würdest, wie großartig...“

„Das tue ich auch, Kjall, wirklich, aber dennoch muss ich darauf bestehen, dass du mir die Sphäre und die Schilde aushändigst.“

„Man sollte wohl nie seine Helden treffen.“

Kjall tat einen mächtigen Satz nach vorne, auf den Takuri-Spiegel zu, der an die Wand gelehnt war. Wenn Kjall durch diesen Spiegel trat, würde an seiner Stelle wohl in Kürze ein verdutzter Troll in dieser Höhle stehen, oder gar ein garstiger Wardrak. Nichts, mit dem ich nicht klarkommen könnte. Dennoch wäre es äußerst ärgerlich, wenn Kjall und seine Sphäre mir durch die Lappen gehen würden. Darum verstärkte ich den Griff um den Sturmschild an meiner rechten Hand und sandte eine gezielte Druckwelle durch die Luft, welche Kjall von den Füßen holte und einen gezackten Spalt in die Spiegeloberfläche schlug. Takuri-Asche stob daraus hervor und glühte leicht auf, während ein Partikel nach dem anderen an einen Ort gesogen wurde – ganz ähnlich, wie es der glühende Tiefsand aus dem Geheimen See mit den Zeitportalen der Sphäre tat.

Wie ich so abgelenkt die Takuri-Asche betrachtete, entging mir völlig, dass Kjall zurück zum Sockel des Feuerschilds hechtete und sich seinen Feuerschild schnappte. Plötzlich wurde mir mulmig zumute. Hoffnung, redete ich mir selbst zu, das wichtigste ist es, die Hoffnung nicht zu verlieren.

„Kjall“, setzte ich an, „Du weißt nicht, was es mit diesem Schild auf sich hat...“

Kjall brüllte irgendetwas Unverständliches und hielt den Feuerschild in die Höhe, über dessen ornamentierte Oberfläche ein unheilvoller schwarzer Glanz huschte. Schwarz-silbernes Feuer brach aus den Wänden des Museums hervor und leckte an den verschiedensten Artefakten, die da standen. Die Flammen fraßen sich in Sekundenschnelle durch den Hadrischen Spiegel, die andorische Flöte, die Überreste des Hadrischen Stundenglases, dutzende Notizen und Artefakte, Holz und Stahl.

Die Hitze schlug mir entgegen und versengte mir den Bart. Ich drehte Kjall den Rücken zu, in der Hoffnung, dass mein eigener Feuerschild das Feuer ein wenig absorbieren würde. Wo waren die Roteisenschilde, wenn man sie brauchte?! Dann griff ich erneut nach der Macht des Sturmschilds und fühlte den vertrauten Druck in meinen Ohren, als ich die Luft auseinanderzog und eine Barriere zwischen dem flammenden Inferno im Museum und der Luftblase, die mich umgab, schuf.

Kjall stand inmitten des wütenden Feuers und sah zu, wie es sich ins Gestein fraß, seine Artefakte verzehrte und seine Arbeit in Funken aufgehen ließ. Er stand einfach da, ein höllisches Lachen seiner Kehle entspringend, während sein Lebenswerk in Flammen aufging. Eine Gänsehaut lief mir über den Rücken. Ich musste dem ein Ende setzen, sofort.

Das silberne Feuer brandete immer noch an die luftleere Barriere, die meine kleine Luftblase umschloss. Ein drittes und letztes Mal griff ich nach der Macht meines Sturmschildes, spürte die Wassertropfen in der Luft, die Spannungen in der Erde und den Druck auf den Stollen. Während die Barriere um mich herum sich auflöste und das silberne Feuer auf mich zustürzte, zog ich die Luft um Kjalls Kopf auseinander und hüllte ihn in eine luftleere Blase. Kjall japste auf, aber er hatte in all seinen Reisen durch den Limbo der Zeitportale gelernt, seinen Atem auszusetzen. Er griff nach einem der wenigen noch intakten Sockeln, schnappte sich einen schwelenden Hadrischen Kompass und drehte mit seiner freien Hand an der Kompassnadel.

Ein leises Kichern entschlüpfte mir. Wie wollte Kjall mit einem kleinen Kompass, den man an jedem halbwegs besiedelten Meeresufer finden konnte, gegen die Macht des Sturmschilds ausrichten? Leider unterschätzte ich die Reserven des Schilds, welche beinahe entladen waren. Schon ertönte ein blechernes Knirschen und der helle Glanz auf der Oberfläche meines Sturmschilds wurde stumpf. Die Blase um Kjalls Kopf löste sich auf, Wind wirbelte wieder umher und frische Luft strömte in Kjalls Kehle. Purer Hass strahlte aus seinen Augen, als er erneute nach seinem Feuerschild griff und eine silberne Flammenwand aus der Decke des Raumes brach.

Varlion das Flammenschwert, dessen Scheide soeben in Asche aufgegangen war, erhob sich wie von Zauberhand in die Luft und flog Kjall geradewegs in die ausgestreckte freie Hand. Nun mischten sich neben dem silbernen Flammen auch orange Schlieren in die Luft. Kjalls Mund verzerrte sich zu einer wahnsinnigen Fratze, welche aus ihrem Schlund tiefschwarze Flammen direkt in meine Richtung spuckte. Der eiserne Rand meines Bruderschilds verformte sich leicht und der Knubbel in der Schildmitte glühte feuerrot auf, als ich den Schild zwischen die tiefschwarzen Flammen und meine Wenigkeit wuchtete und er die Hitze absorbierte. Ich wusste nicht, wie lange er noch durchhalten könnte.

In einem letzten Effort, gegen die schiere Dunkelheit anzukommen, die mir entgegengeschleudert wurde, griff ich im Geiste gleichzeitig nach dem Sternenschild und dem Bruderschild, verwob Fäden der Hoffnung und Fäden der Stärke und ließ das Geflecht auf Kjall zufliegen. Der Sternenschild flackerte kurz. Lange würden auch die Kraftreserven dieses Schildes nicht mehr halten, aber viel Zeit brauchte ich nicht mehr. Die Verbindung zu Kjall war geschaffen, und so griff ich hoffnungsvoll nach Kjalls Stärke, ließ sie durch mich fließen und vertauschte sie dann mit der Stärke der kleinen Feuermaus, welche zwei Seitengänge von uns entfernt ahnungslos durch die Stollen huschte.

Kjall klappte zusammen wie eine Puppe, der man die Fäden durchgeschnitten hatte. Der Feuerschild und Varlion das Flammenschwert fielen aus seinen schwachen Händen und das Flammeninferno versiegte. Seine Sphäre rollte aus seiner Manteltasche. Ich rannte zu Kjall herüber und trat den Feuerschild zur Seite. Mein Sternenschild flackerte ein letztes Mal auf und sein Brummen verstummte, woraufhin Kjalls Kraft zurück in seinen Körper floss, aber ich war schneller und drückte Kjall mit einem Stiefel zu Boden. Fast wie nebenbei steckte ich seine Sphäre ein.

Kjalls Blick blieb noch einige Augenblicke lang von Hass erfüllt, dann kam er wieder zur Besinnung und schlug sich die Hände vor den Kopf. Ich ließ von ihm ab und trat zurück. Kjall drehte seinen Kopf zur Seite und betrachtete die Zerstörung, die der Feuerschild angerichtet hatte. Das Innere des Museums war nun von einer dicken Ascheschicht überzogen. Steinbrocken lagen kreuz und quer über den Raum verteilt, durchmischt mit verkohlten Holzsplittern und verbogenen Metallplatten. Ein trauriger Anblick.

„Ich... ich weiß nicht, was über mich gekommen ist“, wimmerte Kjall zitternd, mehr zu sich selbst als zu mir. Kraftlos griff er nach einem Stück Kohle – den Überresten des Hadrischen Kompasses – welches unter seinen Fingern zerstob. „Nein, nein, das kann nicht sein, das darf nicht sein.“

Ich ignorierte Kjall fürs erste und begutachtete die zerstörten Artefakte. Die meisten konnte man nicht mal mehr erkennen. Varlion brannte immer noch, aber seine Flammen sahen beinahe süß aus im Vergleich zum Inferno, welches hier soeben gewütet hatte. Kjall und seine Kleidung waren unversehrt und ich war mit einigen hässlichen Rötungen weggekommen. Wir hatten beide unglaubliches Glück gehabt.

Kjall schlug mit der Faust auf den Boden und blickte mir wütend entgegen. Ich konnte es ihm nachfühlen, aber ihm nicht helfen. Nicht, dass er meinen Trost überhaupt angenommen hätte. Und so wandte ich mich mit einem mulmigen Gefühl im Magen wieder ab.

Vor mir lagen die vier mächtigen Schilde aus uralter Zeit gleich in doppelter Ausführung. Die linken vier hatte ich zusammen mit der Sphäre in meiner linken Manteltasche erst vor einigen Tagen von meinem zukünftigen Ich überreicht bekommen und in diesen Kampf mitgebracht. Die rechten vier hatte Kjall aus einer Vielzahl von Stellen in der Vergangenheit erbeutet. Nun war es wohl an der Zeit, mich auf die Socken zu machen. Ich musste die linken vier Schilde an ihre rechtmäßige Stelle in der Geschichte zurückbringen. Aber vorher sollte ich lieber noch die rechten vier Schilde und Kjalls Sphäre meinem früheren Ich überbringen gehen.

Ich ergriff meinen großen Sack und packte alle acht Schilde sorgfältig darin ein. Stolz überkam mich, als ich meine Meisterwerke ein letztes Mal betrachtete, auch wenn ich beim Anblick der Feuerschilde leicht erschauderte. Kjalls Sphäre steckte ich lieber in meine Jackentasche als in den Sack. Die rechte, damit sie ja nicht mit der Sphäre in meiner linken Manteltasche verwechselt wurde. Ich schulterte den Sack unter erheblichem Gerumpel und wandte mich zurück an Kjall:

„Ich schätze, dass es für mich an der Zeit ist zu gehen. Es ist bald Weihnachten und ich habe eine Menge Geschenke auszutragen. Lebewohl, Sohn des Xoll. Es war mir eine Freude, dich kennenzulernen, und es wäre mir eine Freude gewesen, dich noch näher kennenzulernen. Aber das Schicksal scheint andere Pläne für uns zu haben. Mein Beileid für deine Sammlung. Sie war wirklich außergewöhnlich, und das Gewissen, sie erreicht zu haben, kann dir niemand mehr wegnehmen.“

Kjall blickte mich an, nun ganz und gar nicht mehr wütend oder schockiert, nicht einmal mehr enttäuscht. Seine Augen hatten stattdessen einen flehenden Eindruck angenommen:

„Kreatok?“, fragte er stockend.

„Hm?“

„Wie steht es um den Treibstoff in der Sphäre? Wäre... wäre es dir möglich, mich auf eine letzte Reise zu schicken?“, Kjall wies schwach auf sein zerstörtes Museum, „Hier hält mich nichts mehr. Ich träumte ohnehin schon jede Nacht von den Urzeiten Caverns, den bemannten Festen, den Drachen...“

Ich hatte ihn wirklich falsch eingeschätzt. Sein ganzes Lebenswerk war in Flammen aufgegangen und es wäre nur allzu verständlich gewesen, wenn er mich als Auslöser dieses Untergangs angesehen und in die Tiefen des Enran verdammt hätte. Aber hier saß er und bat mich darum, ihm einen letzten Wunsch zu erfüllen. Er tat mir leid.

Ich entgegnete: „Bist du sicher, dass das die beste Entscheidung ist? Was ist mit deinen Freunden hier? Radan, seine Schar?“

„Was soll schon mit ihnen sein, Kreatok?“

Kjalls Blick wurde wieder wässerig, und ich fragte mich erneut, ob ich Kjall unterschätzt hatte. Konnte es sein, dass er herausgefunden hatte, dass ich ihm von Josella hatte einen Bären aufbinden lassen, was Caverns Zukunft anging? Hatte er versucht, die drei ihm damals noch fehlenden mächtigen Schilde zunächst nach der Rückkehr der Helden aus Krahd zu erringen? Hatte er erfahren, welche Schildzwerge gemeinsam mit den vier Schilden im Enran versinken würden?

„Mich hält hier nichts mehr. Bitte.“, wiederholte Kjall, diesmal gefasster. Er schnäuzte sich in den rußigen Bart.

Ich blickte ihn nachdenklich an. Ärger anrichten konnte er in dieser Zeitlinie nicht mehr, das wäre mir doch sicherlich schon aufgefallen. Und in den Urzeiten gab es weit und breit keine Sphäre, mit der Kjall zurückkommen könnte.

„Sprichst du denn überhaupt die Alte Sprache?“, fragte ich den immer noch vor mir knienden Zwerg.

Kjall antwortete ohne zu zögern: „Isklja Dugandr!“

Ich nickte beeindruckt. Ich unterschätzte ihn immer noch. Leiste murmelnd zog ich Kjalls Sphäre hervor, klappte die Hauptstreben zur Seite, legte den linken Schalter um und drehte das kleine Zahnrad im Kontrollzentrum zurück, immer und immer weiter, bis es anschlug. Dann stellte ich die Sphäre auf den Boden und vernahm mit Ehrfurcht, wie es darin klickte, ratterte und surrte, flüssiges Feuer mit Tiefsand vermischt wurde und fröhlich blubberte, bis der vertraute blaue Dampf aus seiner Öffnung gepustet wurde und sich im verkohlten Raum ausbreitete. Was für eine meisterhafte Konstruktion diese Sphäre doch war! Ich konnte es immer noch kaum fassen, dass ich es war, der sie erschaffen hatte.

Kjall blickte mich weiterhin unverändert an. Aus seinem Blick sprach Akzeptanz, und ja, selbst Dankbarkeit. Das zerstörte Museum schien vergessen. Dachte er an glorreiche Zwergenfesten und feuerspuckende Drachen, Meisterwerke der Schmiedekunst und die wackere Bezwingung des unberührten Felsgesteins? Dieses romantische Bild würde ihm noch früh genug vergehen.

Ich klopfte ihm auf die Schulter.

„Tief durchatmen, Kjall. Das wird eine lange Reise für dich.“


\begin{center}
    Weiter geht es in \hypref{Die Legende der entzweiten Zeit (2023)} und \hypref{Der Giftzwerg und das Drachenherz (2023)}.
\end{center}








\newpage
\section{Epilog: Bragor erhält den Bruderschild zurück}

\az{Jahr 62}


Das Rietgras wogte sachte im Wind. Eine massige Gestalt wanderte über das Feld. Nur mit einem Lendenschurz aus Schaffell und großen Stiefeln war sie bekleidet. In ihrer einen Hand hielt sie einen Speer. Auf dem Kopf saßen zwei mächtige Hörner. Das Sonnenlicht ließ die prallen Muskeln des Wesens glänzen. Es war ein Tarus.

Eine kleine Gestalt trat aus dem Schatten eines Baumes hervor und näherte sich dem Tarus. Sie hatte einen großen Sack über die Schulter geworfen. In der ein Hand hielt sie eine Art silbernen Kugel, in der zweiten einen runden Schild. Ein grauer Bart war in den Gürtel der Gestalt gesteckt. Es war ein Zwerg.

„Seid herzlich gegrüßt!“, rief der Zwerg fröhlich und hob dann seine Hand, die den Rundschild festhielt, „Ich bin hier, um unser Geschäft rückgängig zu machen! Der Schild, den Ihr mir verkauft hat, funktioniert gar nicht“

Der Tarus runzelte seine Stirn: „Was meint Ihr damit, dass der Schild nicht funktioniere? Und Ihr seid doch gar nicht derselbe Zwerg, dem ich diesen Schild gerade eben verkauft hatte?“

Der Zwerg lachte schallend auf: „Nicht derselbe Zwerg?! Das ist ein guter Trick, den muss ich mir merken!“

Der Tarus blickte den Zwerg weiterhin verwirrt an. Dieser klopfte dem Tarus gegen den Bauch und meinte: „Nun, ist das hier Euer Schild oder nicht?“

„Nun ja, er sieht schon so aus, wie derjenige, den ich...“

„Wunderbar, dass wäre das ja geklärt. Ich bin nicht zufrieden mit dem Schild und will ihn gerne zurückgeben!“

Der Tarus kniff die Augen zusammen: „Nein, Ihr seht wirklich nicht so aus wie der Zwerg, mit dem ich soeben einen Handel abgeschlossen habe. Und auch wenn ich nicht viel vom Handel verstehen mag... wenn Ihr mir den Schild zurückgebt, dann muss ich Euch das Goldstück wiedergeben, und das würde ich gerne behalten.“

Den Zwerg machte eine wegwerfende Handbewegung: „Dann behaltet das Gold halt. Ich will einfach nur den Schild loswerden!“

Der Tarus kratzte sich am Kopf, kniete sich dann aber brav hin und streckte seinen Arm aus, damit der Zwerg ihm den Rundschild wieder anbinden konnte. Sogar ein besonderes silbernes Seil benutzte der Zwerg. Verstehe einer diese Kontinentalbewohner und ihre Gepflogenheiten.

Als der Zwerg sich aufrichtete und bereits wieder weglaufen wollte, streckte der Tarus seine Hand aus und meinte: „Ihr... nein, dieser andere Zwerg, dem ich diesen Schild verkauft habe... mir wurde geraten, mich an den Baum der Lieder oder eine Kräuterhexe zu wenden. Ich war bereits am Baum der Lieder, und die wollten mir nur eine schwere Ziege aufbinden. Sagt, wisst Ihr, wo ich vielleicht eine Kräuterhexe finden könnte, die sich mit der Verbreitung von Sternkraut auskennt?“

Der Zwerg kratzte nachdenklich seinen bauschigen Bart, zeigte dann mit seinem Finger in den Südwesten und sagte: „An Eurer Stelle würde ich mich an die Hütte hinter dem Hügel dort drüben wenden. Ich habe gehört, dass sich dort erst kürzlich eine Hexe hat blicken lassen, und dass dort eine vorzügliche Fischsuppe mit Lauch und Wernersilie aufgetischt wird. Auch wenn Ihr die reine Wernersilie wohl lieber hättet, wenn ich Euch so ansehe. Vielleicht werdet Ihr dort Eure Auskunft erhalten – oder jemanden finden, der Euch den Schild abnimmt.“

Der Zwerge drehte und schraubte an der seltsamen silbernen Kugel in seiner Hand herum und ein großes blaues Portal bildete sich aus hellblauem Dampf. Der Zwerg blickte prüfend hinein. Dann holte er Anlauf und sprang in das Portal, welches sich zu einem kleinen blauen Funken zusammenzog und mit einem leisen Plopp verschwand.

Der Tarus blieb alleine übrig und kratzte sich an den Hörnern. Was hier abging, war eindeutig nicht alltäglich. Er drehte sich um und sah nachdenklich den Hügel an, hinter dem angeblich eine Hütte lag, in welcher gelegentlich eine Hexe einkehrte.

Noch war der Tag jung.

Zeit, dieser Hütte einen Besuch abzustatten.

